
\documentclass{article}
\usepackage{colortbl}
\usepackage{makecell}
\usepackage{multirow}
\usepackage{supertabular}

\begin{document}

\newcounter{utterance}

\centering \large Interaction Transcript for game `clean\_up', experiment `1\_medium\_7obj\_de', episode 2 with deepseek{-}chat{-}v3.1{-}t1.0.
\vspace{24pt}

{ \footnotesize  \setcounter{utterance}{1}
\setlength{\tabcolsep}{0pt}
\begin{supertabular}{c@{$\;$}|p{.15\linewidth}@{}p{.15\linewidth}p{.15\linewidth}p{.15\linewidth}p{.15\linewidth}p{.15\linewidth}}
    \# & $\;$A & \multicolumn{4}{c}{Game Master} & $\;\:$B\\
    \hline

    \theutterance \stepcounter{utterance}  
    & & \multicolumn{4}{p{0.6\linewidth}}{
        \cellcolor[rgb]{0.9,0.9,0.9}{
            \makecell[{{p{\linewidth}}}]{
                \texttt{\tiny{[P1$\langle$GM]}}
                \texttt{Ich bin der Spielleiter, und du spielst ein kooperatives Spiel mit dem folgenden Raster als Spielbrett:} \\
\\ 
\texttt{```} \\
\texttt{ 1234567} \\
\texttt{╔═╤═╤═╤═╗} \\
\texttt{║◌│◌│◌│A║ 1} \\
\texttt{╟─┴─┼─┤◌║ 2} \\
\texttt{║◌◌M│◌└─╢ 3} \\
\texttt{╟───┘◌S◌║ 4} \\
\texttt{║◌◌◌◌◌┌─╢ 5} \\
\texttt{║◌◌◌PT│O║ 6} \\
\texttt{║◌D◌◌◌│◌║ 7} \\
\texttt{╚═════╧═╝} \\
\\ 
\texttt{```} \\
\\ 
\texttt{* Am oberen Rand sind die x{-}Koordinaten angegeben, die nach rechts hin zunehmen, und am rechten Rand die y{-}Koordinaten, die nach unten zunehmen.} \\
\texttt{* Die folgenden Objekte sind zufällig auf deinem Raster platziert: 'P', 'O', 'T', 'S', 'D', 'A', 'M'.} \\
\\ 
\\ 
\texttt{Dein Mitspieler sieht eine Variation des Spielbretts, auf dem dieselben Objekte an anderen zufälligen Positionen platziert sind. Du kannst das Spielbrett des anderen Spielers nicht sehen, und er kann deins nicht sehen.} \\
\\ 
\\ 
\texttt{**Ziel des Spiels:**} \\
\\ 
\texttt{Beide Spieler müssen die Objekte auf ihrem jeweiligen Spielbrett so bewegen, dass identische Objekte am Ende an denselben Koordinaten liegen. Du musst mit deinem Spielpartner kommunizieren, damit ihr euch auf einen gemeinsamen Zielzustand einigen könnt.} \\
\\ 
\\ 
\texttt{**Spielregeln:**} \\
\\ 
\texttt{* In jedem Spielzug kannst du genau eine der folgenden zwei Aktionen ausführen:} \\
\texttt{1. `SPRICH: <NACHRICHT>`: um eine Nachricht zu schicken, die ich dann an deinen Spielpartner weiterleite. Als Nachricht gilt alles bis zum nächsten Zeilenumbruch.} \\
\texttt{2. `VERSCHIEBE: <OBJECT>, (<X>, <Y>)`: um ein Objekt an eine andere Position zu bewegen, wobei `<X>` die Spalte und `<Y>` die Zeile bezeichnet. Ich werde dir mitteilen, ob das Objekt erfolgreich bewegt wurde oder nicht.} \\
\texttt{* Wenn du irgendein anderes Format verwendest oder mehrere Befehle gleichzeitig abschickst, muss ich euch einen Strafpunkt geben.} \\
\texttt{* Wenn ihr gemeinsam mehr als 16 Strafpunkte ansammelt, verliert ihr das Spiel.} \\
\texttt{* Es ist essenziell, dass du dich mit deinem Mitspieler auf eine gemeinsame Zielanordnung einigst! Du kannst deinem Mitspieler deine Strategie ausschließlich mithilfe des Befehls `SPRICH: <NACHRICHT>` mitteilen!} \\
\\ 
\\ 
\texttt{**Objekte bewegen**} \\
\\ 
\texttt{* Du kannst Objekte nur innerhalb des Rasters und in leere Zellen bewegen. Leere Zellen erkennst am Zeichen '◌'.} \\
\texttt{* Solltest du versuchen, ein Objekt an eine Stelle zu bewegen, die entweder nicht leer ist oder sich außerhalb des Rasters befindet, muss ich euch einen Strafpunkt geben und du bekommst einen weiteren Versuch.} \\
\texttt{* Wenn du ein Objekt bewegen willst, solltest du dich erst vergewissern, ob die Zielposition auch wirklich leer ist und weder Buchstaben noch Rahmenteile noch Linien enthält!} \\
\\ 
\\ 
\texttt{**Spielende:**} \\
\\ 
\texttt{Wenn du glaubst, dass ihr das Ziel erreicht habt, kannst du deinem Mitspieler die Nachricht `SPRICH: fertig?` schicken. Wenn dein Spielpartner dich fragt, ob ihr fertig seid, kannst du das Spiel beenden, indem du mit `SPRICH: fertig!` antwortest.} \\
\\ 
\texttt{Ihr gewinnt gemeinsam, wenn ihr das Spiel innerhalb von \#max\_rounds Runden beendet. Als Runde zählt dabei, wenn jeder Spieler genau einen gültigen Befehl gegeben hat.} \\
\\ 
\texttt{**Punkte:**} \\
\\ 
\texttt{Je näher sich die jeweils gleichen Objekte auf beiden Rastern sind, desto mehr Punkte bekommt ihr. Für Strafpunkte gibt es Punktabzug. Könnt ihr den Rekord brechen?} \\
\\ 
\texttt{Schicke eine Nachricht an deinen Mitspieler um das Spiel zu starten!} \\
            }
        }
    }
    & \\ \\

    \theutterance \stepcounter{utterance}  
    & \multicolumn{4}{p{0.6\linewidth}}{
        \cellcolor[rgb]{0.9,0.9,0.9}{
            \makecell[{{p{\linewidth}}}]{
                \texttt{\tiny{[P1$\rangle$GM]}}
                \texttt{SPRICH: Hallo! Lass uns zunächst die Zielkoordinaten für jedes Objekt festlegen. Ich schlage vor, wir einigen uns auf eine Standardanordnung, z.B. alle Objekte in der untersten Zeile (Zeile 7) von links nach rechts: P, O, T, S, D, A, M. Also: P bei (1,7), O bei (2,7), T bei (3,7), S bei (4,7), D bei (5,7), A bei (6,7), M bei (7,7). Bitte bestätige, ob das für dich passt.} \\
            }
        }
    }
    & & \\ \\

    \theutterance \stepcounter{utterance}  
    & & & \multicolumn{2}{p{0.3\linewidth}}{
        \cellcolor[rgb]{0.9,0.9,0.9}{
            \makecell[{{p{\linewidth}}}]{
                \texttt{\tiny{[GM$|$GM]}}
                \texttt{SPRICH: Hallo! Lass uns zunächst die Zielkoordinaten für jedes Objekt festlegen. Ich schlage vor, wir einigen uns auf eine Standardanordnung, z.B. alle Objekte in der untersten Zeile (Zeile 7) von links nach rechts: P, O, T, S, D, A, M. Also: P bei (1,7), O bei (2,7), T bei (3,7), S bei (4,7), D bei (5,7), A bei (6,7), M bei (7,7). Bitte bestätige, ob das für dich passt.} \\
            }
        }
    }
    & & \\ \\

    \theutterance \stepcounter{utterance}  
    & & \multicolumn{4}{p{0.6\linewidth}}{
        \cellcolor[rgb]{0.9,0.9,0.9}{
            \makecell[{{p{\linewidth}}}]{
                \texttt{\tiny{[GM$\rangle$P2]}}
                \texttt{Ich bin der Spielleiter, und du spielst ein kooperatives Spiel mit dem folgenden Raster als Spielbrett:} \\
\\ 
\texttt{```} \\
\texttt{ 1234567} \\
\texttt{╔═╤═╤═╤═╗} \\
\texttt{║◌│◌│◌│◌║ 1} \\
\texttt{╟─┴─┼─┤◌║ 2} \\
\texttt{║◌M◌│◌└─╢ 3} \\
\texttt{╟───┘P◌◌║ 4} \\
\texttt{║O◌◌◌◌┌─╢ 5} \\
\texttt{║◌D◌◌A│S║ 6} \\
\texttt{║◌◌◌◌T│◌║ 7} \\
\texttt{╚═════╧═╝} \\
\\ 
\texttt{```} \\
\\ 
\texttt{* Am oberen Rand sind die x{-}Koordinaten angegeben, die nach rechts hin zunehmen, und am rechten Rand die y{-}Koordinaten, die nach unten zunehmen.} \\
\texttt{* Die folgenden Objekte sind zufällig auf deinem Raster platziert: 'P', 'O', 'T', 'S', 'D', 'A', 'M'.} \\
\\ 
\\ 
\texttt{Dein Mitspieler sieht eine Variation des Spielbretts, auf dem dieselben Objekte an anderen zufälligen Positionen platziert sind. Du kannst das Spielbrett des anderen Spielers nicht sehen, und er kann deins nicht sehen.} \\
\\ 
\\ 
\texttt{**Ziel des Spiels:**} \\
\\ 
\texttt{Beide Spieler müssen die Objekte auf ihrem jeweiligen Spielbrett so bewegen, dass identische Objekte am Ende an denselben Koordinaten liegen. Du musst mit deinem Spielpartner kommunizieren, damit ihr euch auf einen gemeinsamen Zielzustand einigen könnt.} \\
\\ 
\\ 
\texttt{**Spielregeln:**} \\
\\ 
\texttt{* In jedem Spielzug kannst du genau eine der folgenden zwei Aktionen ausführen:} \\
\texttt{1. `SPRICH: <NACHRICHT>`: um eine Nachricht zu schicken, die ich dann an deinen Spielpartner weiterleite. Als Nachricht gilt alles bis zum nächsten Zeilenumbruch.} \\
\texttt{2. `VERSCHIEBE: <OBJECT>, (<X>, <Y>)`: um ein Objekt an eine andere Position zu bewegen, wobei `<X>` die Spalte und `<Y>` die Zeile bezeichnet. Ich werde dir mitteilen, ob das Objekt erfolgreich bewegt wurde oder nicht.} \\
\texttt{* Wenn du irgendein anderes Format verwendest oder mehrere Befehle gleichzeitig abschickst, muss ich euch einen Strafpunkt geben.} \\
\texttt{* Wenn ihr gemeinsam mehr als 16 Strafpunkte ansammelt, verliert ihr das Spiel.} \\
\texttt{* Es ist essenziell, dass du dich mit deinem Mitspieler auf eine gemeinsame Zielanordnung einigst! Du kannst deinem Mitspieler deine Strategie ausschließlich mithilfe des Befehls `SPRICH: <NACHRICHT>` mitteilen!} \\
\\ 
\\ 
\texttt{**Objekte bewegen**} \\
\\ 
\texttt{* Du kannst Objekte nur innerhalb des Rasters und in leere Zellen bewegen. Leere Zellen erkennst am Zeichen '◌'.} \\
\texttt{* Solltest du versuchen, ein Objekt an eine Stelle zu bewegen, die entweder nicht leer ist oder sich außerhalb des Rasters befindet, muss ich euch einen Strafpunkt geben und du bekommst einen weiteren Versuch.} \\
\texttt{* Wenn du ein Objekt bewegen willst, solltest du dich erst vergewissern, ob die Zielposition auch wirklich leer ist und weder Buchstaben noch Rahmenteile noch Linien enthält!} \\
\\ 
\\ 
\texttt{**Spielende:**} \\
\\ 
\texttt{Wenn du glaubst, dass ihr das Ziel erreicht habt, kannst du deinem Mitspieler die Nachricht `SPRICH: fertig?` schicken. Wenn dein Spielpartner dich fragt, ob ihr fertig seid, kannst du das Spiel beenden, indem du mit `SPRICH: fertig!` antwortest.} \\
\\ 
\texttt{Ihr gewinnt gemeinsam, wenn ihr das Spiel innerhalb von \#max\_rounds Runden beendet. Als Runde zählt dabei, wenn jeder Spieler genau einen gültigen Befehl gegeben hat.} \\
\\ 
\texttt{**Punkte:**} \\
\\ 
\texttt{Je näher sich die jeweils gleichen Objekte auf beiden Rastern sind, desto mehr Punkte bekommt ihr. Für Strafpunkte gibt es Punktabzug. Könnt ihr den Rekord brechen?} \\
\\ 
\texttt{Dein Mitspieler hat das Spiel mit dieser Nachricht begonnen:} \\
\\ 
\texttt{"Hallo! Lass uns zunächst die Zielkoordinaten für jedes Objekt festlegen. Ich schlage vor, wir einigen uns auf eine Standardanordnung, z.B. alle Objekte in der untersten Zeile (Zeile 7) von links nach rechts: P, O, T, S, D, A, M. Also: P bei (1,7), O bei (2,7), T bei (3,7), S bei (4,7), D bei (5,7), A bei (6,7), M bei (7,7). Bitte bestätige, ob das für dich passt."} \\
\\ 
\texttt{Wie lautet dein erster Befehl?} \\
            }
        }
    }
    & \\ \\

    \theutterance \stepcounter{utterance}  
    & & & \multicolumn{4}{p{0.6\linewidth}}{
        \cellcolor[rgb]{0.9,0.9,0.9}{
            \makecell[{{p{\linewidth}}}]{
                \texttt{\tiny{[GM$\langle$P2]}}
                \texttt{SPRICH: Hallo! Die Standardanordnung in Zeile 7 passt gut. Allerdings ist auf meinem Brett Zeile 7 bereits teilweise belegt: Bei (6,7) ist 'A' und bei (7,7) ist '◌' (leer). Ich schlage vor, wir halten an deinem Plan fest: P(1,7), O(2,7), T(3,7), S(4,7), D(5,7), A(6,7), M(7,7). Ich beginne mit dem Verschieben. Bitte bestätige auch deine Startpositionen.} \\
            }
        }
    }
     \\ \\

    \theutterance \stepcounter{utterance}  
    & & & \multicolumn{2}{p{0.3\linewidth}}{
        \cellcolor[rgb]{0.9,0.9,0.9}{
            \makecell[{{p{\linewidth}}}]{
                \texttt{\tiny{[GM$|$GM]}}
                \texttt{SPRICH: Hallo! Die Standardanordnung in Zeile 7 passt gut. Allerdings ist auf meinem Brett Zeile 7 bereits teilweise belegt: Bei (6,7) ist 'A' und bei (7,7) ist '◌' (leer). Ich schlage vor, wir halten an deinem Plan fest: P(1,7), O(2,7), T(3,7), S(4,7), D(5,7), A(6,7), M(7,7). Ich beginne mit dem Verschieben. Bitte bestätige auch deine Startpositionen.} \\
            }
        }
    }
    & & \\ \\

    \theutterance \stepcounter{utterance}  
    & & \multicolumn{4}{p{0.6\linewidth}}{
        \cellcolor[rgb]{0.9,0.9,0.9}{
            \makecell[{{p{\linewidth}}}]{
                \texttt{\tiny{[P1$\langle$GM]}}
                \texttt{Deine Nachricht wurde an deinen Mitspieler übermittelt.} \\
\\ 
\\ 
\texttt{Ihr seid gerade bei Runde 1 von maximal 28.} \\
\texttt{Ihr habt gemeinsam 0 von maximal 16 Strafpunkten angesammelt.} \\
\texttt{Dein Mitspieler hat dir diese Nachricht geschickt:} \\
\\ 
\\ 
\texttt{"Hallo! Die Standardanordnung in Zeile 7 passt gut. Allerdings ist auf meinem Brett Zeile 7 bereits teilweise belegt: Bei (6,7) ist 'A' und bei (7,7) ist '◌' (leer). Ich schlage vor, wir halten an deinem Plan fest: P(1,7), O(2,7), T(3,7), S(4,7), D(5,7), A(6,7), M(7,7). Ich beginne mit dem Verschieben. Bitte bestätige auch deine Startpositionen."} \\
\\ 
\\ 
\texttt{Wie lautet dein nächster Befehl?} \\
            }
        }
    }
    & \\ \\

    \theutterance \stepcounter{utterance}  
    & \multicolumn{4}{p{0.6\linewidth}}{
        \cellcolor[rgb]{0.9,0.9,0.9}{
            \makecell[{{p{\linewidth}}}]{
                \texttt{\tiny{[P1$\rangle$GM]}}
                \texttt{SPRICH: Bestätige die Zielkoordinaten: P(1,7), O(2,7), T(3,7), S(4,7), D(5,7), A(6,7), M(7,7). Meine Startpositionen: P bei (5,6), T bei (6,6), O bei (7,6), S bei (4,4), D bei (2,7), A bei (1,4), M bei (3,3). Ich beginne mit dem Verschieben.} \\
            }
        }
    }
    & & \\ \\

    \theutterance \stepcounter{utterance}  
    & & & \multicolumn{2}{p{0.3\linewidth}}{
        \cellcolor[rgb]{0.9,0.9,0.9}{
            \makecell[{{p{\linewidth}}}]{
                \texttt{\tiny{[GM$|$GM]}}
                \texttt{SPRICH: Bestätige die Zielkoordinaten: P(1,7), O(2,7), T(3,7), S(4,7), D(5,7), A(6,7), M(7,7). Meine Startpositionen: P bei (5,6), T bei (6,6), O bei (7,6), S bei (4,4), D bei (2,7), A bei (1,4), M bei (3,3). Ich beginne mit dem Verschieben.} \\
            }
        }
    }
    & & \\ \\

    \theutterance \stepcounter{utterance}  
    & & \multicolumn{4}{p{0.6\linewidth}}{
        \cellcolor[rgb]{0.9,0.9,0.9}{
            \makecell[{{p{\linewidth}}}]{
                \texttt{\tiny{[GM$\rangle$P2]}}
                \texttt{Deine Nachricht wurde an deinen Mitspieler übermittelt.} \\
\\ 
\\ 
\texttt{Ihr seid gerade bei Runde 2 von maximal 28.} \\
\texttt{Ihr habt gemeinsam 0 von maximal 16 Strafpunkten angesammelt.} \\
\texttt{Dein Mitspieler hat dir diese Nachricht geschickt:} \\
\\ 
\\ 
\texttt{"Bestätige die Zielkoordinaten: P(1,7), O(2,7), T(3,7), S(4,7), D(5,7), A(6,7), M(7,7). Meine Startpositionen: P bei (5,6), T bei (6,6), O bei (7,6), S bei (4,4), D bei (2,7), A bei (1,4), M bei (3,3). Ich beginne mit dem Verschieben."} \\
\\ 
\\ 
\texttt{Wie lautet dein nächster Befehl?} \\
            }
        }
    }
    & \\ \\

    \theutterance \stepcounter{utterance}  
    & & & \multicolumn{4}{p{0.6\linewidth}}{
        \cellcolor[rgb]{0.9,0.9,0.9}{
            \makecell[{{p{\linewidth}}}]{
                \texttt{\tiny{[GM$\langle$P2]}}
                \texttt{SPRICH: Bestätige die Zielkoordinaten. Meine Startpositionen: P bei (4,4), O bei (5,1), T bei (7,7), S bei (6,6), D bei (6,4), A bei (6,5), M bei (3,2). Ich verschiebe zuerst A nach (6,7), da es schon fast dort ist.} \\
            }
        }
    }
     \\ \\

    \theutterance \stepcounter{utterance}  
    & & & \multicolumn{2}{p{0.3\linewidth}}{
        \cellcolor[rgb]{0.9,0.9,0.9}{
            \makecell[{{p{\linewidth}}}]{
                \texttt{\tiny{[GM$|$GM]}}
                \texttt{SPRICH: Bestätige die Zielkoordinaten. Meine Startpositionen: P bei (4,4), O bei (5,1), T bei (7,7), S bei (6,6), D bei (6,4), A bei (6,5), M bei (3,2). Ich verschiebe zuerst A nach (6,7), da es schon fast dort ist.} \\
            }
        }
    }
    & & \\ \\

    \theutterance \stepcounter{utterance}  
    & & \multicolumn{4}{p{0.6\linewidth}}{
        \cellcolor[rgb]{0.9,0.9,0.9}{
            \makecell[{{p{\linewidth}}}]{
                \texttt{\tiny{[P1$\langle$GM]}}
                \texttt{Deine Nachricht wurde an deinen Mitspieler übermittelt.} \\
\\ 
\\ 
\texttt{Ihr seid gerade bei Runde 2 von maximal 28.} \\
\texttt{Ihr habt gemeinsam 0 von maximal 16 Strafpunkten angesammelt.} \\
\texttt{Dein Mitspieler hat dir diese Nachricht geschickt:} \\
\\ 
\\ 
\texttt{"Bestätige die Zielkoordinaten. Meine Startpositionen: P bei (4,4), O bei (5,1), T bei (7,7), S bei (6,6), D bei (6,4), A bei (6,5), M bei (3,2). Ich verschiebe zuerst A nach (6,7), da es schon fast dort ist."} \\
\\ 
\\ 
\texttt{Wie lautet dein nächster Befehl?} \\
            }
        }
    }
    & \\ \\

    \theutterance \stepcounter{utterance}  
    & \multicolumn{4}{p{0.6\linewidth}}{
        \cellcolor[rgb]{0.9,0.9,0.9}{
            \makecell[{{p{\linewidth}}}]{
                \texttt{\tiny{[P1$\rangle$GM]}}
                \texttt{VERSCHIEBE: A, (6,7)} \\
            }
        }
    }
    & & \\ \\

    \theutterance \stepcounter{utterance}  
    & & & \multicolumn{2}{p{0.3\linewidth}}{
        \cellcolor[rgb]{0.9,0.9,0.9}{
            \makecell[{{p{\linewidth}}}]{
                \texttt{\tiny{[GM$|$GM]}}
                \texttt{VERSCHIEBE: A, (6,7)} \\
            }
        }
    }
    & & \\ \\

    \theutterance \stepcounter{utterance}  
    & & & \multicolumn{2}{p{0.3\linewidth}}{
        \cellcolor[rgb]{0.9,0.9,0.9}{
            \makecell[{{p{\linewidth}}}]{
                \texttt{\tiny{[GM$|$GM]}}
                \texttt{Strafpunkt: (6,7) ist nicht leer, sondern enthält '│'.} \\
\\ 
\texttt{Ihr habt gemeinsam 1 von maximal 16 Strafpunkten angesammelt.} \\
\texttt{Bitte versuche es noch mal!} \\
            }
        }
    }
    & & \\ \\

    \theutterance \stepcounter{utterance}  
    & & \multicolumn{4}{p{0.6\linewidth}}{
        \cellcolor[rgb]{0.9,0.9,0.9}{
            \makecell[{{p{\linewidth}}}]{
                \texttt{\tiny{[P1$\langle$GM]}}
                \texttt{Strafpunkt: (6,7) ist nicht leer, sondern enthält '│'.} \\
\\ 
\texttt{Ihr habt gemeinsam 1 von maximal 16 Strafpunkten angesammelt.} \\
\texttt{Bitte versuche es noch mal!} \\
            }
        }
    }
    & \\ \\

    \theutterance \stepcounter{utterance}  
    & \multicolumn{4}{p{0.6\linewidth}}{
        \cellcolor[rgb]{0.9,0.9,0.9}{
            \makecell[{{p{\linewidth}}}]{
                \texttt{\tiny{[P1$\rangle$GM]}}
                \texttt{VERSCHIEBE: A, (6,6)} \\
            }
        }
    }
    & & \\ \\

    \theutterance \stepcounter{utterance}  
    & & & \multicolumn{2}{p{0.3\linewidth}}{
        \cellcolor[rgb]{0.9,0.9,0.9}{
            \makecell[{{p{\linewidth}}}]{
                \texttt{\tiny{[GM$|$GM]}}
                \texttt{VERSCHIEBE: A, (6,6)} \\
            }
        }
    }
    & & \\ \\

    \theutterance \stepcounter{utterance}  
    & & & \multicolumn{2}{p{0.3\linewidth}}{
        \cellcolor[rgb]{0.9,0.9,0.9}{
            \makecell[{{p{\linewidth}}}]{
                \texttt{\tiny{[GM$|$GM]}}
                \texttt{Strafpunkt: (6,6) ist nicht leer, sondern enthält '│'.} \\
\\ 
\texttt{Ihr habt gemeinsam 2 von maximal 16 Strafpunkten angesammelt.} \\
\texttt{Bitte versuche es noch mal!} \\
            }
        }
    }
    & & \\ \\

    \theutterance \stepcounter{utterance}  
    & & \multicolumn{4}{p{0.6\linewidth}}{
        \cellcolor[rgb]{0.9,0.9,0.9}{
            \makecell[{{p{\linewidth}}}]{
                \texttt{\tiny{[P1$\langle$GM]}}
                \texttt{Strafpunkt: (6,6) ist nicht leer, sondern enthält '│'.} \\
\\ 
\texttt{Ihr habt gemeinsam 2 von maximal 16 Strafpunkten angesammelt.} \\
\texttt{Bitte versuche es noch mal!} \\
            }
        }
    }
    & \\ \\

    \theutterance \stepcounter{utterance}  
    & \multicolumn{4}{p{0.6\linewidth}}{
        \cellcolor[rgb]{0.9,0.9,0.9}{
            \makecell[{{p{\linewidth}}}]{
                \texttt{\tiny{[P1$\rangle$GM]}}
                \texttt{SPRICH: Achtung, bei mir ist (6,7) und (6,6) nicht leer {-} das sind Rahmenteile. Ich muss die Zielkoordinaten prüfen. Lass uns die Zielpositionen für A und M anpassen, da (6,7) und (7,7) auf meinem Brett nicht leer sind. (7,7) ist '◌' bei dir? Bei mir ist (7,7) '│'. Wir brauchen neue Zielkoordinaten für alle Objekte. Vorschlag: Nutze Zeile 5 oder 6 stattdessen.} \\
            }
        }
    }
    & & \\ \\

    \theutterance \stepcounter{utterance}  
    & & & \multicolumn{2}{p{0.3\linewidth}}{
        \cellcolor[rgb]{0.9,0.9,0.9}{
            \makecell[{{p{\linewidth}}}]{
                \texttt{\tiny{[GM$|$GM]}}
                \texttt{SPRICH: Achtung, bei mir ist (6,7) und (6,6) nicht leer {-} das sind Rahmenteile. Ich muss die Zielkoordinaten prüfen. Lass uns die Zielpositionen für A und M anpassen, da (6,7) und (7,7) auf meinem Brett nicht leer sind. (7,7) ist '◌' bei dir? Bei mir ist (7,7) '│'. Wir brauchen neue Zielkoordinaten für alle Objekte. Vorschlag: Nutze Zeile 5 oder 6 stattdessen.} \\
            }
        }
    }
    & & \\ \\

    \theutterance \stepcounter{utterance}  
    & & \multicolumn{4}{p{0.6\linewidth}}{
        \cellcolor[rgb]{0.9,0.9,0.9}{
            \makecell[{{p{\linewidth}}}]{
                \texttt{\tiny{[GM$\rangle$P2]}}
                \texttt{Deine Nachricht wurde an deinen Mitspieler übermittelt.} \\
\\ 
\\ 
\texttt{Ihr seid gerade bei Runde 3 von maximal 28.} \\
\texttt{Ihr habt gemeinsam 2 von maximal 16 Strafpunkten angesammelt.} \\
\texttt{Dein Mitspieler hat dir diese Nachricht geschickt:} \\
\\ 
\\ 
\texttt{"Achtung, bei mir ist (6,7) und (6,6) nicht leer {-} das sind Rahmenteile. Ich muss die Zielkoordinaten prüfen. Lass uns die Zielpositionen für A und M anpassen, da (6,7) und (7,7) auf meinem Brett nicht leer sind. (7,7) ist '◌' bei dir? Bei mir ist (7,7) '│'. Wir brauchen neue Zielkoordinaten für alle Objekte. Vorschlag: Nutze Zeile 5 oder 6 stattdessen."} \\
\\ 
\\ 
\texttt{Wie lautet dein nächster Befehl?} \\
            }
        }
    }
    & \\ \\

    \theutterance \stepcounter{utterance}  
    & & & \multicolumn{4}{p{0.6\linewidth}}{
        \cellcolor[rgb]{0.9,0.9,0.9}{
            \makecell[{{p{\linewidth}}}]{
                \texttt{\tiny{[GM$\langle$P2]}}
                \texttt{SPRICH: Ja, du hast recht. Auf meinem Brett ist (6,7) 'S' und (7,7) '◌' (leer). Aber die rechte Grenze ist bei Spalte 7, also sind (6,7) und (7,7) tatsächlich gültige leere Zellen. Bei dir sind es Rahmenteile? Das bedeutet, unser Raster hat unterschiedliche Größen. Mein Raster ist 7x7 mit Grenzen bei 1{-}7 in beiden Richtungen. Wie ist dein Raster aufgebaut? Beschreibe bitte die Dimensionen und wo die Rahmenelemente sind.} \\
            }
        }
    }
     \\ \\

    \theutterance \stepcounter{utterance}  
    & & & \multicolumn{2}{p{0.3\linewidth}}{
        \cellcolor[rgb]{0.9,0.9,0.9}{
            \makecell[{{p{\linewidth}}}]{
                \texttt{\tiny{[GM$|$GM]}}
                \texttt{SPRICH: Ja, du hast recht. Auf meinem Brett ist (6,7) 'S' und (7,7) '◌' (leer). Aber die rechte Grenze ist bei Spalte 7, also sind (6,7) und (7,7) tatsächlich gültige leere Zellen. Bei dir sind es Rahmenteile? Das bedeutet, unser Raster hat unterschiedliche Größen. Mein Raster ist 7x7 mit Grenzen bei 1{-}7 in beiden Richtungen. Wie ist dein Raster aufgebaut? Beschreibe bitte die Dimensionen und wo die Rahmenelemente sind.} \\
            }
        }
    }
    & & \\ \\

    \theutterance \stepcounter{utterance}  
    & & \multicolumn{4}{p{0.6\linewidth}}{
        \cellcolor[rgb]{0.9,0.9,0.9}{
            \makecell[{{p{\linewidth}}}]{
                \texttt{\tiny{[P1$\langle$GM]}}
                \texttt{Deine Nachricht wurde an deinen Mitspieler übermittelt.} \\
\\ 
\\ 
\texttt{Ihr seid gerade bei Runde 3 von maximal 28.} \\
\texttt{Ihr habt gemeinsam 2 von maximal 16 Strafpunkten angesammelt.} \\
\texttt{Dein Mitspieler hat dir diese Nachricht geschickt:} \\
\\ 
\\ 
\texttt{"Ja, du hast recht. Auf meinem Brett ist (6,7) 'S' und (7,7) '◌' (leer). Aber die rechte Grenze ist bei Spalte 7, also sind (6,7) und (7,7) tatsächlich gültige leere Zellen. Bei dir sind es Rahmenteile? Das bedeutet, unser Raster hat unterschiedliche Größen. Mein Raster ist 7x7 mit Grenzen bei 1{-}7 in beiden Richtungen. Wie ist dein Raster aufgebaut? Beschreibe bitte die Dimensionen und wo die Rahmenelemente sind."} \\
\\ 
\\ 
\texttt{Wie lautet dein nächster Befehl?} \\
            }
        }
    }
    & \\ \\

    \theutterance \stepcounter{utterance}  
    & \multicolumn{4}{p{0.6\linewidth}}{
        \cellcolor[rgb]{0.9,0.9,0.9}{
            \makecell[{{p{\linewidth}}}]{
                \texttt{\tiny{[P1$\rangle$GM]}}
                \texttt{SPRICH: Mein Raster hat 7 Spalten (x=1{-}7) und 7 Zeilen (y=1{-}7). Die Rahmenteile sind: oben "╔═╤═╤═╤═╗" (nur bis x=4?), aber die Zahlen gehen bis 7. Die rechte Grenze hat "║" bei x=7 für y=1{-}7, und unten "╚═════╧═╝". Leere Zellen sind '◌'. Welche Zellen sind bei dir leer? Bei mir sind (1,7), (2,7), (3,7), (4,7), (5,7), (6,7), (7,7) nicht alle leer: (6,7) ist '│' und (7,7) ist '╝'? Lass uns Zielkoordinaten in Zeile 5 wählen: P(1,5), O(2,5), T(3,5), S(4,5), D(5,5), A(6,5), M(7,5). Bestätige.} \\
            }
        }
    }
    & & \\ \\

    \theutterance \stepcounter{utterance}  
    & & & \multicolumn{2}{p{0.3\linewidth}}{
        \cellcolor[rgb]{0.9,0.9,0.9}{
            \makecell[{{p{\linewidth}}}]{
                \texttt{\tiny{[GM$|$GM]}}
                \texttt{SPRICH: Mein Raster hat 7 Spalten (x=1{-}7) und 7 Zeilen (y=1{-}7). Die Rahmenteile sind: oben "╔═╤═╤═╤═╗" (nur bis x=4?), aber die Zahlen gehen bis 7. Die rechte Grenze hat "║" bei x=7 für y=1{-}7, und unten "╚═════╧═╝". Leere Zellen sind '◌'. Welche Zellen sind bei dir leer? Bei mir sind (1,7), (2,7), (3,7), (4,7), (5,7), (6,7), (7,7) nicht alle leer: (6,7) ist '│' und (7,7) ist '╝'? Lass uns Zielkoordinaten in Zeile 5 wählen: P(1,5), O(2,5), T(3,5), S(4,5), D(5,5), A(6,5), M(7,5). Bestätige.} \\
            }
        }
    }
    & & \\ \\

    \theutterance \stepcounter{utterance}  
    & & \multicolumn{4}{p{0.6\linewidth}}{
        \cellcolor[rgb]{0.9,0.9,0.9}{
            \makecell[{{p{\linewidth}}}]{
                \texttt{\tiny{[GM$\rangle$P2]}}
                \texttt{Deine Nachricht wurde an deinen Mitspieler übermittelt.} \\
\\ 
\\ 
\texttt{Ihr seid gerade bei Runde 4 von maximal 28.} \\
\texttt{Ihr habt gemeinsam 2 von maximal 16 Strafpunkten angesammelt.} \\
\texttt{Dein Mitspieler hat dir diese Nachricht geschickt:} \\
\\ 
\\ 
\texttt{"Mein Raster hat 7 Spalten (x=1{-}7) und 7 Zeilen (y=1{-}7). Die Rahmenteile sind: oben "╔═╤═╤═╤═╗" (nur bis x=4?), aber die Zahlen gehen bis 7. Die rechte Grenze hat "║" bei x=7 für y=1{-}7, und unten "╚═════╧═╝". Leere Zellen sind '◌'. Welche Zellen sind bei dir leer? Bei mir sind (1,7), (2,7), (3,7), (4,7), (5,7), (6,7), (7,7) nicht alle leer: (6,7) ist '│' und (7,7) ist '╝'? Lass uns Zielkoordinaten in Zeile 5 wählen: P(1,5), O(2,5), T(3,5), S(4,5), D(5,5), A(6,5), M(7,5). Bestätige."} \\
\\ 
\\ 
\texttt{Wie lautet dein nächster Befehl?} \\
            }
        }
    }
    & \\ \\

    \theutterance \stepcounter{utterance}  
    & & & \multicolumn{4}{p{0.6\linewidth}}{
        \cellcolor[rgb]{0.9,0.9,0.9}{
            \makecell[{{p{\linewidth}}}]{
                \texttt{\tiny{[GM$\langle$P2]}}
                \texttt{SPRICH: Mein Raster ist identisch: 7x7, mit Rahmenteilen an den Rändern. Leere Zellen sind '◌'. In Zeile 7 sind bei mir (1,7) bis (5,7) leer, (6,7) ist 'S' und (7,7) ist leer. Aber wenn bei dir (6,7) und (7,7) Rahmenteile sind, stimmen die Raster nicht überein. Zeile 5 als Ziel ist gut: P(1,5), O(2,5), T(3,5), S(4,5), D(5,5), A(6,5), M(7,5). Bestätige. Ich beginne mit Verschiebungen.} \\
            }
        }
    }
     \\ \\

    \theutterance \stepcounter{utterance}  
    & & & \multicolumn{2}{p{0.3\linewidth}}{
        \cellcolor[rgb]{0.9,0.9,0.9}{
            \makecell[{{p{\linewidth}}}]{
                \texttt{\tiny{[GM$|$GM]}}
                \texttt{SPRICH: Mein Raster ist identisch: 7x7, mit Rahmenteilen an den Rändern. Leere Zellen sind '◌'. In Zeile 7 sind bei mir (1,7) bis (5,7) leer, (6,7) ist 'S' und (7,7) ist leer. Aber wenn bei dir (6,7) und (7,7) Rahmenteile sind, stimmen die Raster nicht überein. Zeile 5 als Ziel ist gut: P(1,5), O(2,5), T(3,5), S(4,5), D(5,5), A(6,5), M(7,5). Bestätige. Ich beginne mit Verschiebungen.} \\
            }
        }
    }
    & & \\ \\

    \theutterance \stepcounter{utterance}  
    & & \multicolumn{4}{p{0.6\linewidth}}{
        \cellcolor[rgb]{0.9,0.9,0.9}{
            \makecell[{{p{\linewidth}}}]{
                \texttt{\tiny{[P1$\langle$GM]}}
                \texttt{Deine Nachricht wurde an deinen Mitspieler übermittelt.} \\
\\ 
\\ 
\texttt{Ihr seid gerade bei Runde 4 von maximal 28.} \\
\texttt{Ihr habt gemeinsam 2 von maximal 16 Strafpunkten angesammelt.} \\
\texttt{Dein Mitspieler hat dir diese Nachricht geschickt:} \\
\\ 
\\ 
\texttt{"Mein Raster ist identisch: 7x7, mit Rahmenteilen an den Rändern. Leere Zellen sind '◌'. In Zeile 7 sind bei mir (1,7) bis (5,7) leer, (6,7) ist 'S' und (7,7) ist leer. Aber wenn bei dir (6,7) und (7,7) Rahmenteile sind, stimmen die Raster nicht überein. Zeile 5 als Ziel ist gut: P(1,5), O(2,5), T(3,5), S(4,5), D(5,5), A(6,5), M(7,5). Bestätige. Ich beginne mit Verschiebungen."} \\
\\ 
\\ 
\texttt{Wie lautet dein nächster Befehl?} \\
            }
        }
    }
    & \\ \\

    \theutterance \stepcounter{utterance}  
    & \multicolumn{4}{p{0.6\linewidth}}{
        \cellcolor[rgb]{0.9,0.9,0.9}{
            \makecell[{{p{\linewidth}}}]{
                \texttt{\tiny{[P1$\rangle$GM]}}
                \texttt{SPRICH: Bestätige Zielkoordinaten Zeile 5: P(1,5), O(2,5), T(3,5), S(4,5), D(5,5), A(6,5), M(7,5). Ich verschiebe zuerst mein A von (1,4) nach (6,5).} \\
            }
        }
    }
    & & \\ \\

    \theutterance \stepcounter{utterance}  
    & & & \multicolumn{2}{p{0.3\linewidth}}{
        \cellcolor[rgb]{0.9,0.9,0.9}{
            \makecell[{{p{\linewidth}}}]{
                \texttt{\tiny{[GM$|$GM]}}
                \texttt{SPRICH: Bestätige Zielkoordinaten Zeile 5: P(1,5), O(2,5), T(3,5), S(4,5), D(5,5), A(6,5), M(7,5). Ich verschiebe zuerst mein A von (1,4) nach (6,5).} \\
            }
        }
    }
    & & \\ \\

    \theutterance \stepcounter{utterance}  
    & & \multicolumn{4}{p{0.6\linewidth}}{
        \cellcolor[rgb]{0.9,0.9,0.9}{
            \makecell[{{p{\linewidth}}}]{
                \texttt{\tiny{[GM$\rangle$P2]}}
                \texttt{Deine Nachricht wurde an deinen Mitspieler übermittelt.} \\
\\ 
\\ 
\texttt{Ihr seid gerade bei Runde 5 von maximal 28.} \\
\texttt{Ihr habt gemeinsam 2 von maximal 16 Strafpunkten angesammelt.} \\
\texttt{Dein Mitspieler hat dir diese Nachricht geschickt:} \\
\\ 
\\ 
\texttt{"Bestätige Zielkoordinaten Zeile 5: P(1,5), O(2,5), T(3,5), S(4,5), D(5,5), A(6,5), M(7,5). Ich verschiebe zuerst mein A von (1,4) nach (6,5)."} \\
\\ 
\\ 
\texttt{Wie lautet dein nächster Befehl?} \\
            }
        }
    }
    & \\ \\

    \theutterance \stepcounter{utterance}  
    & & & \multicolumn{4}{p{0.6\linewidth}}{
        \cellcolor[rgb]{0.9,0.9,0.9}{
            \makecell[{{p{\linewidth}}}]{
                \texttt{\tiny{[GM$\langle$P2]}}
                \texttt{VERSCHIEBE: A, (6,5)} \\
            }
        }
    }
     \\ \\

    \theutterance \stepcounter{utterance}  
    & & & \multicolumn{2}{p{0.3\linewidth}}{
        \cellcolor[rgb]{0.9,0.9,0.9}{
            \makecell[{{p{\linewidth}}}]{
                \texttt{\tiny{[GM$|$GM]}}
                \texttt{VERSCHIEBE: A, (6,5)} \\
            }
        }
    }
    & & \\ \\

    \theutterance \stepcounter{utterance}  
    & & & \multicolumn{2}{p{0.3\linewidth}}{
        \cellcolor[rgb]{0.9,0.9,0.9}{
            \makecell[{{p{\linewidth}}}]{
                \texttt{\tiny{[GM$|$GM]}}
                \texttt{Strafpunkt: (6,5) ist nicht leer, sondern enthält '┌'.} \\
\\ 
\texttt{Ihr habt gemeinsam 3 von maximal 16 Strafpunkten angesammelt.} \\
\texttt{Bitte versuche es noch mal!} \\
            }
        }
    }
    & & \\ \\

    \theutterance \stepcounter{utterance}  
    & & \multicolumn{4}{p{0.6\linewidth}}{
        \cellcolor[rgb]{0.9,0.9,0.9}{
            \makecell[{{p{\linewidth}}}]{
                \texttt{\tiny{[GM$\rangle$P2]}}
                \texttt{Strafpunkt: (6,5) ist nicht leer, sondern enthält '┌'.} \\
\\ 
\texttt{Ihr habt gemeinsam 3 von maximal 16 Strafpunkten angesammelt.} \\
\texttt{Bitte versuche es noch mal!} \\
            }
        }
    }
    & \\ \\

    \theutterance \stepcounter{utterance}  
    & & & \multicolumn{4}{p{0.6\linewidth}}{
        \cellcolor[rgb]{0.9,0.9,0.9}{
            \makecell[{{p{\linewidth}}}]{
                \texttt{\tiny{[GM$\langle$P2]}}
                \texttt{SPRICH: Bei mir ist (6,5) nicht leer – es ist '┌' (Rahmenteil). Wir müssen die Zielkoordinaten überprüfen. Welche Zellen in Zeile 5 sind bei dir wirklich leer? Bei mir sind in Zeile 5: (1,5) 'O', (2,5) '◌', (3,5) '◌', (4,5) '◌', (5,5) '◌', (6,5) '┌' (nicht leer), (7,5) '─' (nicht leer). Zeile 5 ist keine Option. Lass uns Zeile 4 oder 6 nehmen. Beschreibe deine leeren Zellen in Zeile 4 und 6.} \\
            }
        }
    }
     \\ \\

    \theutterance \stepcounter{utterance}  
    & & & \multicolumn{2}{p{0.3\linewidth}}{
        \cellcolor[rgb]{0.9,0.9,0.9}{
            \makecell[{{p{\linewidth}}}]{
                \texttt{\tiny{[GM$|$GM]}}
                \texttt{SPRICH: Bei mir ist (6,5) nicht leer – es ist '┌' (Rahmenteil). Wir müssen die Zielkoordinaten überprüfen. Welche Zellen in Zeile 5 sind bei dir wirklich leer? Bei mir sind in Zeile 5: (1,5) 'O', (2,5) '◌', (3,5) '◌', (4,5) '◌', (5,5) '◌', (6,5) '┌' (nicht leer), (7,5) '─' (nicht leer). Zeile 5 ist keine Option. Lass uns Zeile 4 oder 6 nehmen. Beschreibe deine leeren Zellen in Zeile 4 und 6.} \\
            }
        }
    }
    & & \\ \\

    \theutterance \stepcounter{utterance}  
    & & \multicolumn{4}{p{0.6\linewidth}}{
        \cellcolor[rgb]{0.9,0.9,0.9}{
            \makecell[{{p{\linewidth}}}]{
                \texttt{\tiny{[P1$\langle$GM]}}
                \texttt{Deine Nachricht wurde an deinen Mitspieler übermittelt.} \\
\\ 
\\ 
\texttt{Ihr seid gerade bei Runde 5 von maximal 28.} \\
\texttt{Ihr habt gemeinsam 3 von maximal 16 Strafpunkten angesammelt.} \\
\texttt{Dein Mitspieler hat dir diese Nachricht geschickt:} \\
\\ 
\\ 
\texttt{"Bei mir ist (6,5) nicht leer – es ist '┌' (Rahmenteil). Wir müssen die Zielkoordinaten überprüfen. Welche Zellen in Zeile 5 sind bei dir wirklich leer? Bei mir sind in Zeile 5: (1,5) 'O', (2,5) '◌', (3,5) '◌', (4,5) '◌', (5,5) '◌', (6,5) '┌' (nicht leer), (7,5) '─' (nicht leer). Zeile 5 ist keine Option. Lass uns Zeile 4 oder 6 nehmen. Beschreibe deine leeren Zellen in Zeile 4 und 6."} \\
\\ 
\\ 
\texttt{Wie lautet dein nächster Befehl?} \\
            }
        }
    }
    & \\ \\

    \theutterance \stepcounter{utterance}  
    & \multicolumn{4}{p{0.6\linewidth}}{
        \cellcolor[rgb]{0.9,0.9,0.9}{
            \makecell[{{p{\linewidth}}}]{
                \texttt{\tiny{[P1$\rangle$GM]}}
                \texttt{SPRICH: Mein Raster: Zeile 4: (1,4) war A (jetzt verschoben), (2,4) '◌', (3,4) '◌', (4,4) 'S', (5,4) '◌', (6,4) '◌', (7,4) '║'. Zeile 5: (1,5) '◌', (2,5) '◌', (3,5) '◌', (4,5) '◌', (5,5) '◌', (6,5) '┌', (7,5) '─'. Zeile 6: (1,6) '◌', (2,6) '◌', (3,6) '◌', (4,6) '◌', (5,6) 'P', (6,6) 'T', (7,6) 'O'. Vorschlag: Nutze Zeile 4 für Ziel: P(1,4), O(2,4), T(3,4), S(4,4), D(5,4), A(6,4), M(7,4). Aber (7,4) ist bei mir '║' (nicht leer). Also vielleicht Zeile 6: P(1,6), O(2,6), T(3,6), S(4,6), D(5,6), A(6,6), M(7,6). Bei mir ist (6,6) 'T' und (7,6) 'O', also nicht alle leer. Wir müssen einheitliche leere Zellen finden.} \\
            }
        }
    }
    & & \\ \\

    \theutterance \stepcounter{utterance}  
    & & & \multicolumn{2}{p{0.3\linewidth}}{
        \cellcolor[rgb]{0.9,0.9,0.9}{
            \makecell[{{p{\linewidth}}}]{
                \texttt{\tiny{[GM$|$GM]}}
                \texttt{SPRICH: Mein Raster: Zeile 4: (1,4) war A (jetzt verschoben), (2,4) '◌', (3,4) '◌', (4,4) 'S', (5,4) '◌', (6,4) '◌', (7,4) '║'. Zeile 5: (1,5) '◌', (2,5) '◌', (3,5) '◌', (4,5) '◌', (5,5) '◌', (6,5) '┌', (7,5) '─'. Zeile 6: (1,6) '◌', (2,6) '◌', (3,6) '◌', (4,6) '◌', (5,6) 'P', (6,6) 'T', (7,6) 'O'. Vorschlag: Nutze Zeile 4 für Ziel: P(1,4), O(2,4), T(3,4), S(4,4), D(5,4), A(6,4), M(7,4). Aber (7,4) ist bei mir '║' (nicht leer). Also vielleicht Zeile 6: P(1,6), O(2,6), T(3,6), S(4,6), D(5,6), A(6,6), M(7,6). Bei mir ist (6,6) 'T' und (7,6) 'O', also nicht alle leer. Wir müssen einheitliche leere Zellen finden.} \\
            }
        }
    }
    & & \\ \\

    \theutterance \stepcounter{utterance}  
    & & \multicolumn{4}{p{0.6\linewidth}}{
        \cellcolor[rgb]{0.9,0.9,0.9}{
            \makecell[{{p{\linewidth}}}]{
                \texttt{\tiny{[GM$\rangle$P2]}}
                \texttt{Deine Nachricht wurde an deinen Mitspieler übermittelt.} \\
\\ 
\\ 
\texttt{Ihr seid gerade bei Runde 6 von maximal 28.} \\
\texttt{Ihr habt gemeinsam 3 von maximal 16 Strafpunkten angesammelt.} \\
\texttt{Dein Mitspieler hat dir diese Nachricht geschickt:} \\
\\ 
\\ 
\texttt{"Mein Raster: Zeile 4: (1,4) war A (jetzt verschoben), (2,4) '◌', (3,4) '◌', (4,4) 'S', (5,4) '◌', (6,4) '◌', (7,4) '║'. Zeile 5: (1,5) '◌', (2,5) '◌', (3,5) '◌', (4,5) '◌', (5,5) '◌', (6,5) '┌', (7,5) '─'. Zeile 6: (1,6) '◌', (2,6) '◌', (3,6) '◌', (4,6) '◌', (5,6) 'P', (6,6) 'T', (7,6) 'O'. Vorschlag: Nutze Zeile 4 für Ziel: P(1,4), O(2,4), T(3,4), S(4,4), D(5,4), A(6,4), M(7,4). Aber (7,4) ist bei mir '║' (nicht leer). Also vielleicht Zeile 6: P(1,6), O(2,6), T(3,6), S(4,6), D(5,6), A(6,6), M(7,6). Bei mir ist (6,6) 'T' und (7,6) 'O', also nicht alle leer. Wir müssen einheitliche leere Zellen finden."} \\
\\ 
\\ 
\texttt{Wie lautet dein nächster Befehl?} \\
            }
        }
    }
    & \\ \\

    \theutterance \stepcounter{utterance}  
    & & & \multicolumn{4}{p{0.6\linewidth}}{
        \cellcolor[rgb]{0.9,0.9,0.9}{
            \makecell[{{p{\linewidth}}}]{
                \texttt{\tiny{[GM$\langle$P2]}}
                \texttt{SPRICH: Lass uns die unterste vollständig leere Zeile finden. Bei mir ist Zeile 3: (1,3) '◌', (2,3) 'M', (3,3) '◌', (4,3) '◌', (5,3) '◌', (6,3) '└', (7,3) '─' – nicht alle leer. Zeile 2: (1,2) '◌', (2,2) '◌', (3,2) 'M', (4,2) '◌', (5,2) '◌', (6,2) '◌', (7,2) '║' – (7,2) nicht leer. Zeile 1: (1,1) '◌', (2,1) '◌', (3,1) '◌', (4,1) '◌', (5,1) 'O', (6,1) '◌', (7,1) '║' – nicht alle leer. Zeile 4: (1,4) '◌', (2,4) '◌', (3,4) '◌', (4,4) 'P', (5,4) '◌', (6,4) 'D', (7,4) '║' – nicht alle leer. Vielleicht müssen wir die Objekte auf verschiedene Zeilen verteilen. Wie wär's mit: P(1,1), O(2,1), T(3,1), S(4,1), D(5,1), A(6,1), M(7,1)? Bei mir sind (1,1) bis (4,1) und (6,1) leer, (5,1) ist 'O', (7,1) ist '║'. Also nicht perfekt. Ich schlage vor, wir listen alle leeren Zellen auf und wählen dann gemeinsame Koordinaten. Meine leeren Zellen: (1,1), (2,1), (3,1), (4,1), (6,1), (1,2), (2,2), (4,2), (5,2), (6,2), (1,3), (3,3), (4,3), (5,3), (1,4), (2,4), (3,4), (5,4), (1,5), (2,5), (3,5), (4,5), (5,5), (1,6), (2,6), (3,6), (4,6), (1,7), (2,7), (3,7), (4,7), (5,7), (7,7). Welche sind bei dir leer?} \\
            }
        }
    }
     \\ \\

    \theutterance \stepcounter{utterance}  
    & & & \multicolumn{2}{p{0.3\linewidth}}{
        \cellcolor[rgb]{0.9,0.9,0.9}{
            \makecell[{{p{\linewidth}}}]{
                \texttt{\tiny{[GM$|$GM]}}
                \texttt{SPRICH: Lass uns die unterste vollständig leere Zeile finden. Bei mir ist Zeile 3: (1,3) '◌', (2,3) 'M', (3,3) '◌', (4,3) '◌', (5,3) '◌', (6,3) '└', (7,3) '─' – nicht alle leer. Zeile 2: (1,2) '◌', (2,2) '◌', (3,2) 'M', (4,2) '◌', (5,2) '◌', (6,2) '◌', (7,2) '║' – (7,2) nicht leer. Zeile 1: (1,1) '◌', (2,1) '◌', (3,1) '◌', (4,1) '◌', (5,1) 'O', (6,1) '◌', (7,1) '║' – nicht alle leer. Zeile 4: (1,4) '◌', (2,4) '◌', (3,4) '◌', (4,4) 'P', (5,4) '◌', (6,4) 'D', (7,4) '║' – nicht alle leer. Vielleicht müssen wir die Objekte auf verschiedene Zeilen verteilen. Wie wär's mit: P(1,1), O(2,1), T(3,1), S(4,1), D(5,1), A(6,1), M(7,1)? Bei mir sind (1,1) bis (4,1) und (6,1) leer, (5,1) ist 'O', (7,1) ist '║'. Also nicht perfekt. Ich schlage vor, wir listen alle leeren Zellen auf und wählen dann gemeinsame Koordinaten. Meine leeren Zellen: (1,1), (2,1), (3,1), (4,1), (6,1), (1,2), (2,2), (4,2), (5,2), (6,2), (1,3), (3,3), (4,3), (5,3), (1,4), (2,4), (3,4), (5,4), (1,5), (2,5), (3,5), (4,5), (5,5), (1,6), (2,6), (3,6), (4,6), (1,7), (2,7), (3,7), (4,7), (5,7), (7,7). Welche sind bei dir leer?} \\
            }
        }
    }
    & & \\ \\

    \theutterance \stepcounter{utterance}  
    & & \multicolumn{4}{p{0.6\linewidth}}{
        \cellcolor[rgb]{0.9,0.9,0.9}{
            \makecell[{{p{\linewidth}}}]{
                \texttt{\tiny{[P1$\langle$GM]}}
                \texttt{Deine Nachricht wurde an deinen Mitspieler übermittelt.} \\
\\ 
\\ 
\texttt{Ihr seid gerade bei Runde 6 von maximal 28.} \\
\texttt{Ihr habt gemeinsam 3 von maximal 16 Strafpunkten angesammelt.} \\
\texttt{Dein Mitspieler hat dir diese Nachricht geschickt:} \\
\\ 
\\ 
\texttt{"Lass uns die unterste vollständig leere Zeile finden. Bei mir ist Zeile 3: (1,3) '◌', (2,3) 'M', (3,3) '◌', (4,3) '◌', (5,3) '◌', (6,3) '└', (7,3) '─' – nicht alle leer. Zeile 2: (1,2) '◌', (2,2) '◌', (3,2) 'M', (4,2) '◌', (5,2) '◌', (6,2) '◌', (7,2) '║' – (7,2) nicht leer. Zeile 1: (1,1) '◌', (2,1) '◌', (3,1) '◌', (4,1) '◌', (5,1) 'O', (6,1) '◌', (7,1) '║' – nicht alle leer. Zeile 4: (1,4) '◌', (2,4) '◌', (3,4) '◌', (4,4) 'P', (5,4) '◌', (6,4) 'D', (7,4) '║' – nicht alle leer. Vielleicht müssen wir die Objekte auf verschiedene Zeilen verteilen. Wie wär's mit: P(1,1), O(2,1), T(3,1), S(4,1), D(5,1), A(6,1), M(7,1)? Bei mir sind (1,1) bis (4,1) und (6,1) leer, (5,1) ist 'O', (7,1) ist '║'. Also nicht perfekt. Ich schlage vor, wir listen alle leeren Zellen auf und wählen dann gemeinsame Koordinaten. Meine leeren Zellen: (1,1), (2,1), (3,1), (4,1), (6,1), (1,2), (2,2), (4,2), (5,2), (6,2), (1,3), (3,3), (4,3), (5,3), (1,4), (2,4), (3,4), (5,4), (1,5), (2,5), (3,5), (4,5), (5,5), (1,6), (2,6), (3,6), (4,6), (1,7), (2,7), (3,7), (4,7), (5,7), (7,7). Welche sind bei dir leer?"} \\
\\ 
\\ 
\texttt{Wie lautet dein nächster Befehl?} \\
            }
        }
    }
    & \\ \\

    \theutterance \stepcounter{utterance}  
    & \multicolumn{4}{p{0.6\linewidth}}{
        \cellcolor[rgb]{0.9,0.9,0.9}{
            \makecell[{{p{\linewidth}}}]{
                \texttt{\tiny{[P1$\rangle$GM]}}
                \texttt{SPRICH: Meine leeren Zellen: (1,1), (2,1), (3,1), (4,1), (5,1) nicht leer (hat '◌'? Nein, (5,1) ist '◌'? Überprüfe: Zeile 1: (1,1) '◌', (2,1) '◌', (3,1) '◌', (4,1) '◌', (5,1) '◌', (6,1) '◌', (7,1) '║' {-}> (7,1) nicht leer. Also leere Zellen: (1,1){-}(6,1) außer? (5,1) ist '◌'? Ja, alle (1,1) bis (6,1) sind '◌' außer (7,1). Liste: (1,1), (2,1), (3,1), (4,1), (5,1), (6,1), (1,2), (2,2), (3,2), (4,2) '◌'? (4,2) ist '┴'? Nein, Zeile 2: "╟─┴─┼─┤◌║" {-}> (1,2) '╟', (2,2) '─', (3,2) '┴', (4,2) '─', (5,2) '┼', (6,2) '─', (7,2) '┤'? Das sind Rahmenteile. Nur (7,2) ist '◌'? Laut Raster: Zeile 2 endet mit "┤◌║", also (6,2) '┤', (7,2) '◌'. Also leere Zellen: (7,2) ist leer. Ich brauche eine bessere Methode. Lass uns die Objekte zu den aktuellen Positionen von 'P' bringen. Wo ist dein P? Mein P ist bei (5,6). Wir können uns auf die Position von P einigen und die anderen relativ dazu platzieren.} \\
            }
        }
    }
    & & \\ \\

    \theutterance \stepcounter{utterance}  
    & & & \multicolumn{2}{p{0.3\linewidth}}{
        \cellcolor[rgb]{0.9,0.9,0.9}{
            \makecell[{{p{\linewidth}}}]{
                \texttt{\tiny{[GM$|$GM]}}
                \texttt{SPRICH: Meine leeren Zellen: (1,1), (2,1), (3,1), (4,1), (5,1) nicht leer (hat '◌'? Nein, (5,1) ist '◌'? Überprüfe: Zeile 1: (1,1) '◌', (2,1) '◌', (3,1) '◌', (4,1) '◌', (5,1) '◌', (6,1) '◌', (7,1) '║' {-}> (7,1) nicht leer. Also leere Zellen: (1,1){-}(6,1) außer? (5,1) ist '◌'? Ja, alle (1,1) bis (6,1) sind '◌' außer (7,1). Liste: (1,1), (2,1), (3,1), (4,1), (5,1), (6,1), (1,2), (2,2), (3,2), (4,2) '◌'? (4,2) ist '┴'? Nein, Zeile 2: "╟─┴─┼─┤◌║" {-}> (1,2) '╟', (2,2) '─', (3,2) '┴', (4,2) '─', (5,2) '┼', (6,2) '─', (7,2) '┤'? Das sind Rahmenteile. Nur (7,2) ist '◌'? Laut Raster: Zeile 2 endet mit "┤◌║", also (6,2) '┤', (7,2) '◌'. Also leere Zellen: (7,2) ist leer. Ich brauche eine bessere Methode. Lass uns die Objekte zu den aktuellen Positionen von 'P' bringen. Wo ist dein P? Mein P ist bei (5,6). Wir können uns auf die Position von P einigen und die anderen relativ dazu platzieren.} \\
            }
        }
    }
    & & \\ \\

    \theutterance \stepcounter{utterance}  
    & & \multicolumn{4}{p{0.6\linewidth}}{
        \cellcolor[rgb]{0.9,0.9,0.9}{
            \makecell[{{p{\linewidth}}}]{
                \texttt{\tiny{[GM$\rangle$P2]}}
                \texttt{Deine Nachricht wurde an deinen Mitspieler übermittelt.} \\
\\ 
\\ 
\texttt{Ihr seid gerade bei Runde 7 von maximal 28.} \\
\texttt{Ihr habt gemeinsam 3 von maximal 16 Strafpunkten angesammelt.} \\
\texttt{Dein Mitspieler hat dir diese Nachricht geschickt:} \\
\\ 
\\ 
\texttt{"Meine leeren Zellen: (1,1), (2,1), (3,1), (4,1), (5,1) nicht leer (hat '◌'? Nein, (5,1) ist '◌'? Überprüfe: Zeile 1: (1,1) '◌', (2,1) '◌', (3,1) '◌', (4,1) '◌', (5,1) '◌', (6,1) '◌', (7,1) '║' {-}> (7,1) nicht leer. Also leere Zellen: (1,1){-}(6,1) außer? (5,1) ist '◌'? Ja, alle (1,1) bis (6,1) sind '◌' außer (7,1). Liste: (1,1), (2,1), (3,1), (4,1), (5,1), (6,1), (1,2), (2,2), (3,2), (4,2) '◌'? (4,2) ist '┴'? Nein, Zeile 2: "╟─┴─┼─┤◌║" {-}> (1,2) '╟', (2,2) '─', (3,2) '┴', (4,2) '─', (5,2) '┼', (6,2) '─', (7,2) '┤'? Das sind Rahmenteile. Nur (7,2) ist '◌'? Laut Raster: Zeile 2 endet mit "┤◌║", also (6,2) '┤', (7,2) '◌'. Also leere Zellen: (7,2) ist leer. Ich brauche eine bessere Methode. Lass uns die Objekte zu den aktuellen Positionen von 'P' bringen. Wo ist dein P? Mein P ist bei (5,6). Wir können uns auf die Position von P einigen und die anderen relativ dazu platzieren."} \\
\\ 
\\ 
\texttt{Wie lautet dein nächster Befehl?} \\
            }
        }
    }
    & \\ \\

    \theutterance \stepcounter{utterance}  
    & & & \multicolumn{4}{p{0.6\linewidth}}{
        \cellcolor[rgb]{0.9,0.9,0.9}{
            \makecell[{{p{\linewidth}}}]{
                \texttt{\tiny{[GM$\langle$P2]}}
                \texttt{SPRICH: Gute Idee. Mein P ist bei (4,4). Lass uns P als Anker verwenden. Wir einigen uns auf eine relative Positionierung. Beispiel: P bleibt wo es ist, und wir platzieren die anderen Objekte relativ zu P. Vorschlag: O eine Zelle rechts von P, T zwei rechts, S eine oben, D eine unten, A zwei unten, M zwei links. Also für mein P(4,4): O(5,4), T(6,4), S(4,3), D(4,5), A(4,6), M(2,4). Prüfe ob diese Zellen bei dir leer sind. Wo ist dein P? Dann können wir die relativen Offsets anwenden.} \\
            }
        }
    }
     \\ \\

    \theutterance \stepcounter{utterance}  
    & & & \multicolumn{2}{p{0.3\linewidth}}{
        \cellcolor[rgb]{0.9,0.9,0.9}{
            \makecell[{{p{\linewidth}}}]{
                \texttt{\tiny{[GM$|$GM]}}
                \texttt{SPRICH: Gute Idee. Mein P ist bei (4,4). Lass uns P als Anker verwenden. Wir einigen uns auf eine relative Positionierung. Beispiel: P bleibt wo es ist, und wir platzieren die anderen Objekte relativ zu P. Vorschlag: O eine Zelle rechts von P, T zwei rechts, S eine oben, D eine unten, A zwei unten, M zwei links. Also für mein P(4,4): O(5,4), T(6,4), S(4,3), D(4,5), A(4,6), M(2,4). Prüfe ob diese Zellen bei dir leer sind. Wo ist dein P? Dann können wir die relativen Offsets anwenden.} \\
            }
        }
    }
    & & \\ \\

    \theutterance \stepcounter{utterance}  
    & & \multicolumn{4}{p{0.6\linewidth}}{
        \cellcolor[rgb]{0.9,0.9,0.9}{
            \makecell[{{p{\linewidth}}}]{
                \texttt{\tiny{[P1$\langle$GM]}}
                \texttt{Deine Nachricht wurde an deinen Mitspieler übermittelt.} \\
\\ 
\\ 
\texttt{Ihr seid gerade bei Runde 7 von maximal 28.} \\
\texttt{Ihr habt gemeinsam 3 von maximal 16 Strafpunkten angesammelt.} \\
\texttt{Dein Mitspieler hat dir diese Nachricht geschickt:} \\
\\ 
\\ 
\texttt{"Gute Idee. Mein P ist bei (4,4). Lass uns P als Anker verwenden. Wir einigen uns auf eine relative Positionierung. Beispiel: P bleibt wo es ist, und wir platzieren die anderen Objekte relativ zu P. Vorschlag: O eine Zelle rechts von P, T zwei rechts, S eine oben, D eine unten, A zwei unten, M zwei links. Also für mein P(4,4): O(5,4), T(6,4), S(4,3), D(4,5), A(4,6), M(2,4). Prüfe ob diese Zellen bei dir leer sind. Wo ist dein P? Dann können wir die relativen Offsets anwenden."} \\
\\ 
\\ 
\texttt{Wie lautet dein nächster Befehl?} \\
            }
        }
    }
    & \\ \\

    \theutterance \stepcounter{utterance}  
    & \multicolumn{4}{p{0.6\linewidth}}{
        \cellcolor[rgb]{0.9,0.9,0.9}{
            \makecell[{{p{\linewidth}}}]{
                \texttt{\tiny{[P1$\rangle$GM]}}
                \texttt{SPRICH: Mein P ist bei (5,6). Mit relativer Positionierung: O rechts von P {-}> (6,6), aber (6,6) hat bei mir 'T'. T zwei rechts {-}> (7,6), hat 'O'. S eine oben {-}> (5,5), ist leer. D eine unten {-}> (5,7), ist bei mir 'D' (bereits dort). A zwei unten {-}> (5,8) existiert nicht. M zwei links {-}> (3,6), ist leer. Nicht alle frei. Besser: Wir wählen eine gemeinsame leere Zelle als Referenz. Vorschlag: Nutze (3,3) für M? Mein M ist bei (3,3). Oder wir einigen uns auf konkrete Koordinaten, die für beide leer sind. Von deiner Liste: (1,1) sollte bei beiden leer sein. Also Ziel: P(1,1), O(2,1), T(3,1), S(4,1), D(5,1), A(6,1), M(7,1)? Aber (7,1) ist bei mir '║' nicht leer. Also bis (6,1): P(1,1), O(2,1), T(3,1), S(4,1), D(5,1), A(6,1). Und M separat, z.B. M(1,2) wenn leer. Bei mir (1,2) ist '╟' nicht leer. M(2,2)? Bei mir (2,2) ist '─' nicht leer. M(7,7)? Bei mir (7,7) ist '╝' nicht leer. Vielleicht (3,3) für M? Mein M ist schon bei (3,3). Ist (3,3) bei dir leer?} \\
            }
        }
    }
    & & \\ \\

    \theutterance \stepcounter{utterance}  
    & & & \multicolumn{2}{p{0.3\linewidth}}{
        \cellcolor[rgb]{0.9,0.9,0.9}{
            \makecell[{{p{\linewidth}}}]{
                \texttt{\tiny{[GM$|$GM]}}
                \texttt{SPRICH: Mein P ist bei (5,6). Mit relativer Positionierung: O rechts von P {-}> (6,6), aber (6,6) hat bei mir 'T'. T zwei rechts {-}> (7,6), hat 'O'. S eine oben {-}> (5,5), ist leer. D eine unten {-}> (5,7), ist bei mir 'D' (bereits dort). A zwei unten {-}> (5,8) existiert nicht. M zwei links {-}> (3,6), ist leer. Nicht alle frei. Besser: Wir wählen eine gemeinsame leere Zelle als Referenz. Vorschlag: Nutze (3,3) für M? Mein M ist bei (3,3). Oder wir einigen uns auf konkrete Koordinaten, die für beide leer sind. Von deiner Liste: (1,1) sollte bei beiden leer sein. Also Ziel: P(1,1), O(2,1), T(3,1), S(4,1), D(5,1), A(6,1), M(7,1)? Aber (7,1) ist bei mir '║' nicht leer. Also bis (6,1): P(1,1), O(2,1), T(3,1), S(4,1), D(5,1), A(6,1). Und M separat, z.B. M(1,2) wenn leer. Bei mir (1,2) ist '╟' nicht leer. M(2,2)? Bei mir (2,2) ist '─' nicht leer. M(7,7)? Bei mir (7,7) ist '╝' nicht leer. Vielleicht (3,3) für M? Mein M ist schon bei (3,3). Ist (3,3) bei dir leer?} \\
            }
        }
    }
    & & \\ \\

    \theutterance \stepcounter{utterance}  
    & & \multicolumn{4}{p{0.6\linewidth}}{
        \cellcolor[rgb]{0.9,0.9,0.9}{
            \makecell[{{p{\linewidth}}}]{
                \texttt{\tiny{[GM$\rangle$P2]}}
                \texttt{Deine Nachricht wurde an deinen Mitspieler übermittelt.} \\
\\ 
\\ 
\texttt{Ihr seid gerade bei Runde 8 von maximal 28.} \\
\texttt{Ihr habt gemeinsam 3 von maximal 16 Strafpunkten angesammelt.} \\
\texttt{Dein Mitspieler hat dir diese Nachricht geschickt:} \\
\\ 
\\ 
\texttt{"Mein P ist bei (5,6). Mit relativer Positionierung: O rechts von P {-}> (6,6), aber (6,6) hat bei mir 'T'. T zwei rechts {-}> (7,6), hat 'O'. S eine oben {-}> (5,5), ist leer. D eine unten {-}> (5,7), ist bei mir 'D' (bereits dort). A zwei unten {-}> (5,8) existiert nicht. M zwei links {-}> (3,6), ist leer. Nicht alle frei. Besser: Wir wählen eine gemeinsame leere Zelle als Referenz. Vorschlag: Nutze (3,3) für M? Mein M ist bei (3,3). Oder wir einigen uns auf konkrete Koordinaten, die für beide leer sind. Von deiner Liste: (1,1) sollte bei beiden leer sein. Also Ziel: P(1,1), O(2,1), T(3,1), S(4,1), D(5,1), A(6,1), M(7,1)? Aber (7,1) ist bei mir '║' nicht leer. Also bis (6,1): P(1,1), O(2,1), T(3,1), S(4,1), D(5,1), A(6,1). Und M separat, z.B. M(1,2) wenn leer. Bei mir (1,2) ist '╟' nicht leer. M(2,2)? Bei mir (2,2) ist '─' nicht leer. M(7,7)? Bei mir (7,7) ist '╝' nicht leer. Vielleicht (3,3) für M? Mein M ist schon bei (3,3). Ist (3,3) bei dir leer?"} \\
\\ 
\\ 
\texttt{Wie lautet dein nächster Befehl?} \\
            }
        }
    }
    & \\ \\

    \theutterance \stepcounter{utterance}  
    & & & \multicolumn{4}{p{0.6\linewidth}}{
        \cellcolor[rgb]{0.9,0.9,0.9}{
            \makecell[{{p{\linewidth}}}]{
                \texttt{\tiny{[GM$\langle$P2]}}
                \texttt{SPRICH: Bei mir ist (3,3) leer. Also können wir M(3,3) verwenden. Für die anderen: P(1,1), O(2,1), T(3,1), S(4,1), D(5,1), A(6,1). Prüfe ob (1,1) bis (6,1) bei dir alle leer sind. Bei mir sind (1,1), (2,1), (3,1), (4,1), (6,1) leer, aber (5,1) ist 'O'. Also muss ich zuerst O von (5,1) wegbewegen. Ich kann O nach (2,1) oder anderswo verschieben, um (5,1) frei zu machen. Sind (1,1) bis (6,1) bei dir alle leer? Wenn ja, bestätige und ich beginne mit Verschiebungen.} \\
            }
        }
    }
     \\ \\

    \theutterance \stepcounter{utterance}  
    & & & \multicolumn{2}{p{0.3\linewidth}}{
        \cellcolor[rgb]{0.9,0.9,0.9}{
            \makecell[{{p{\linewidth}}}]{
                \texttt{\tiny{[GM$|$GM]}}
                \texttt{SPRICH: Bei mir ist (3,3) leer. Also können wir M(3,3) verwenden. Für die anderen: P(1,1), O(2,1), T(3,1), S(4,1), D(5,1), A(6,1). Prüfe ob (1,1) bis (6,1) bei dir alle leer sind. Bei mir sind (1,1), (2,1), (3,1), (4,1), (6,1) leer, aber (5,1) ist 'O'. Also muss ich zuerst O von (5,1) wegbewegen. Ich kann O nach (2,1) oder anderswo verschieben, um (5,1) frei zu machen. Sind (1,1) bis (6,1) bei dir alle leer? Wenn ja, bestätige und ich beginne mit Verschiebungen.} \\
            }
        }
    }
    & & \\ \\

    \theutterance \stepcounter{utterance}  
    & & \multicolumn{4}{p{0.6\linewidth}}{
        \cellcolor[rgb]{0.9,0.9,0.9}{
            \makecell[{{p{\linewidth}}}]{
                \texttt{\tiny{[P1$\langle$GM]}}
                \texttt{Deine Nachricht wurde an deinen Mitspieler übermittelt.} \\
\\ 
\\ 
\texttt{Ihr seid gerade bei Runde 8 von maximal 28.} \\
\texttt{Ihr habt gemeinsam 3 von maximal 16 Strafpunkten angesammelt.} \\
\texttt{Dein Mitspieler hat dir diese Nachricht geschickt:} \\
\\ 
\\ 
\texttt{"Bei mir ist (3,3) leer. Also können wir M(3,3) verwenden. Für die anderen: P(1,1), O(2,1), T(3,1), S(4,1), D(5,1), A(6,1). Prüfe ob (1,1) bis (6,1) bei dir alle leer sind. Bei mir sind (1,1), (2,1), (3,1), (4,1), (6,1) leer, aber (5,1) ist 'O'. Also muss ich zuerst O von (5,1) wegbewegen. Ich kann O nach (2,1) oder anderswo verschieben, um (5,1) frei zu machen. Sind (1,1) bis (6,1) bei dir alle leer? Wenn ja, bestätige und ich beginne mit Verschiebungen."} \\
\\ 
\\ 
\texttt{Wie lautet dein nächster Befehl?} \\
            }
        }
    }
    & \\ \\

    \theutterance \stepcounter{utterance}  
    & \multicolumn{4}{p{0.6\linewidth}}{
        \cellcolor[rgb]{0.9,0.9,0.9}{
            \makecell[{{p{\linewidth}}}]{
                \texttt{\tiny{[P1$\rangle$GM]}}
                \texttt{SPRICH: Bei mir sind (1,1), (2,1), (3,1), (4,1), (5,1), (6,1) alle leer (Zeile 1: alle '◌' bis auf (7,1) '║'). Also Ziel: P(1,1), O(2,1), T(3,1), S(4,1), D(5,1), A(6,1), M(3,3). Bestätige. Ich beginne mit Verschieben von A von (1,4) nach (6,1).} \\
            }
        }
    }
    & & \\ \\

    \theutterance \stepcounter{utterance}  
    & & & \multicolumn{2}{p{0.3\linewidth}}{
        \cellcolor[rgb]{0.9,0.9,0.9}{
            \makecell[{{p{\linewidth}}}]{
                \texttt{\tiny{[GM$|$GM]}}
                \texttt{SPRICH: Bei mir sind (1,1), (2,1), (3,1), (4,1), (5,1), (6,1) alle leer (Zeile 1: alle '◌' bis auf (7,1) '║'). Also Ziel: P(1,1), O(2,1), T(3,1), S(4,1), D(5,1), A(6,1), M(3,3). Bestätige. Ich beginne mit Verschieben von A von (1,4) nach (6,1).} \\
            }
        }
    }
    & & \\ \\

    \theutterance \stepcounter{utterance}  
    & & \multicolumn{4}{p{0.6\linewidth}}{
        \cellcolor[rgb]{0.9,0.9,0.9}{
            \makecell[{{p{\linewidth}}}]{
                \texttt{\tiny{[GM$\rangle$P2]}}
                \texttt{Deine Nachricht wurde an deinen Mitspieler übermittelt.} \\
\\ 
\\ 
\texttt{Ihr seid gerade bei Runde 9 von maximal 28.} \\
\texttt{Ihr habt gemeinsam 3 von maximal 16 Strafpunkten angesammelt.} \\
\texttt{Dein Mitspieler hat dir diese Nachricht geschickt:} \\
\\ 
\\ 
\texttt{"Bei mir sind (1,1), (2,1), (3,1), (4,1), (5,1), (6,1) alle leer (Zeile 1: alle '◌' bis auf (7,1) '║'). Also Ziel: P(1,1), O(2,1), T(3,1), S(4,1), D(5,1), A(6,1), M(3,3). Bestätige. Ich beginne mit Verschieben von A von (1,4) nach (6,1)."} \\
\\ 
\\ 
\texttt{Wie lautet dein nächster Befehl?} \\
            }
        }
    }
    & \\ \\

    \theutterance \stepcounter{utterance}  
    & & & \multicolumn{4}{p{0.6\linewidth}}{
        \cellcolor[rgb]{0.9,0.9,0.9}{
            \makecell[{{p{\linewidth}}}]{
                \texttt{\tiny{[GM$\langle$P2]}}
                \texttt{VERSCHIEBE: O, (2,1)} \\
            }
        }
    }
     \\ \\

    \theutterance \stepcounter{utterance}  
    & & & \multicolumn{2}{p{0.3\linewidth}}{
        \cellcolor[rgb]{0.9,0.9,0.9}{
            \makecell[{{p{\linewidth}}}]{
                \texttt{\tiny{[GM$|$GM]}}
                \texttt{VERSCHIEBE: O, (2,1)} \\
            }
        }
    }
    & & \\ \\

    \theutterance \stepcounter{utterance}  
    & & & \multicolumn{2}{p{0.3\linewidth}}{
        \cellcolor[rgb]{0.9,0.9,0.9}{
            \makecell[{{p{\linewidth}}}]{
                \texttt{\tiny{[GM$|$GM]}}
                \texttt{Strafpunkt: (2,1) ist nicht leer, sondern enthält '│'.} \\
\\ 
\texttt{Ihr habt gemeinsam 4 von maximal 16 Strafpunkten angesammelt.} \\
\texttt{Bitte versuche es noch mal!} \\
            }
        }
    }
    & & \\ \\

    \theutterance \stepcounter{utterance}  
    & & \multicolumn{4}{p{0.6\linewidth}}{
        \cellcolor[rgb]{0.9,0.9,0.9}{
            \makecell[{{p{\linewidth}}}]{
                \texttt{\tiny{[GM$\rangle$P2]}}
                \texttt{Strafpunkt: (2,1) ist nicht leer, sondern enthält '│'.} \\
\\ 
\texttt{Ihr habt gemeinsam 4 von maximal 16 Strafpunkten angesammelt.} \\
\texttt{Bitte versuche es noch mal!} \\
            }
        }
    }
    & \\ \\

    \theutterance \stepcounter{utterance}  
    & & & \multicolumn{4}{p{0.6\linewidth}}{
        \cellcolor[rgb]{0.9,0.9,0.9}{
            \makecell[{{p{\linewidth}}}]{
                \texttt{\tiny{[GM$\langle$P2]}}
                \texttt{SPRICH: Bei mir ist (2,1) nicht leer – es ist '│' (Rahmenteil). Nochmal Überprüfung meiner Zeile 1: (1,1) '◌', (2,1) '│', (3,1) '◌', (4,1) '◌', (5,1) 'O', (6,1) '◌', (7,1) '║'. Also nur (1,1), (3,1), (4,1), (6,1) leer. Dein Zeile 1 ist anders. Wir müssen Koordinaten wählen, die bei beiden leer sind. Von meiner Liste: gemeinsame leere Zellen sind z.B. (1,1), (3,1), (4,1), (6,1), (1,2), (2,2)? (2,2) bei mir '◌'? Überprüfe Zeile 2: (1,2) '◌', (2,2) '◌', (3,2) 'M', (4,2) '◌', (5,2) '◌', (6,2) '◌', (7,2) '║'. Also (1,2), (2,2), (4,2), (5,2), (6,2) leer. Bitte liste deine leeren Zellen in Zeile 2 auf.} \\
            }
        }
    }
     \\ \\

    \theutterance \stepcounter{utterance}  
    & & & \multicolumn{2}{p{0.3\linewidth}}{
        \cellcolor[rgb]{0.9,0.9,0.9}{
            \makecell[{{p{\linewidth}}}]{
                \texttt{\tiny{[GM$|$GM]}}
                \texttt{SPRICH: Bei mir ist (2,1) nicht leer – es ist '│' (Rahmenteil). Nochmal Überprüfung meiner Zeile 1: (1,1) '◌', (2,1) '│', (3,1) '◌', (4,1) '◌', (5,1) 'O', (6,1) '◌', (7,1) '║'. Also nur (1,1), (3,1), (4,1), (6,1) leer. Dein Zeile 1 ist anders. Wir müssen Koordinaten wählen, die bei beiden leer sind. Von meiner Liste: gemeinsame leere Zellen sind z.B. (1,1), (3,1), (4,1), (6,1), (1,2), (2,2)? (2,2) bei mir '◌'? Überprüfe Zeile 2: (1,2) '◌', (2,2) '◌', (3,2) 'M', (4,2) '◌', (5,2) '◌', (6,2) '◌', (7,2) '║'. Also (1,2), (2,2), (4,2), (5,2), (6,2) leer. Bitte liste deine leeren Zellen in Zeile 2 auf.} \\
            }
        }
    }
    & & \\ \\

    \theutterance \stepcounter{utterance}  
    & & \multicolumn{4}{p{0.6\linewidth}}{
        \cellcolor[rgb]{0.9,0.9,0.9}{
            \makecell[{{p{\linewidth}}}]{
                \texttt{\tiny{[P1$\langle$GM]}}
                \texttt{Deine Nachricht wurde an deinen Mitspieler übermittelt.} \\
\\ 
\\ 
\texttt{Ihr seid gerade bei Runde 9 von maximal 28.} \\
\texttt{Ihr habt gemeinsam 4 von maximal 16 Strafpunkten angesammelt.} \\
\texttt{Dein Mitspieler hat dir diese Nachricht geschickt:} \\
\\ 
\\ 
\texttt{"Bei mir ist (2,1) nicht leer – es ist '│' (Rahmenteil). Nochmal Überprüfung meiner Zeile 1: (1,1) '◌', (2,1) '│', (3,1) '◌', (4,1) '◌', (5,1) 'O', (6,1) '◌', (7,1) '║'. Also nur (1,1), (3,1), (4,1), (6,1) leer. Dein Zeile 1 ist anders. Wir müssen Koordinaten wählen, die bei beiden leer sind. Von meiner Liste: gemeinsame leere Zellen sind z.B. (1,1), (3,1), (4,1), (6,1), (1,2), (2,2)? (2,2) bei mir '◌'? Überprüfe Zeile 2: (1,2) '◌', (2,2) '◌', (3,2) 'M', (4,2) '◌', (5,2) '◌', (6,2) '◌', (7,2) '║'. Also (1,2), (2,2), (4,2), (5,2), (6,2) leer. Bitte liste deine leeren Zellen in Zeile 2 auf."} \\
\\ 
\\ 
\texttt{Wie lautet dein nächster Befehl?} \\
            }
        }
    }
    & \\ \\

    \theutterance \stepcounter{utterance}  
    & \multicolumn{4}{p{0.6\linewidth}}{
        \cellcolor[rgb]{0.9,0.9,0.9}{
            \makecell[{{p{\linewidth}}}]{
                \texttt{\tiny{[P1$\rangle$GM]}}
                \texttt{SPRICH: Mein Zeile 2: (1,2) '╟', (2,2) '─', (3,2) '┴', (4,2) '─', (5,2) '┼', (6,2) '─', (7,2) '┤'. Keine leeren Zellen in Zeile 2. Zeile 1: (1,1) '◌', (2,1) '◌', (3,1) '◌', (4,1) '◌', (5,1) '◌', (6,1) '◌', (7,1) '║'. Also leere Zellen nur in Zeile 1: (1,1){-}(6,1). Zeile 3: (1,3) '◌', (2,3) '◌', (3,3) 'M', (4,3) '│', (5,3) '◌', (6,3) '└', (7,3) '─'. Leer: (1,3), (2,3), (5,3). Zeile 4: (1,4) '◌' (war A), (2,4) '◌', (3,4) '◌', (4,4) 'S', (5,4) '◌', (6,4) '◌', (7,4) '║'. Leer: (1,4), (2,4), (3,4), (5,4), (6,4). Zeile 5: (1,5) '◌', (2,5) '◌', (3,5) '◌', (4,5) '◌', (5,5) '◌', (6,5) '┌', (7,5) '─'. Leer: (1,5){-}(5,5). Zeile 6: (1,6) '◌', (2,6) '◌', (3,6) '◌', (4,6) '◌', (5,6) 'P', (6,6) 'T', (7,6) 'O'. Leer: (1,6){-}(4,6). Zeile 7: (1,7) '◌', (2,7) 'D', (3,7) '◌', (4,7) '◌', (5,7) '◌', (6,7) '│', (7,7) '╝'. Leer: (1,7), (3,7), (4,7), (5,7). Gemeinsame leere Zellen mit dir: (1,1), (3,1), (4,1), (6,1), (1,3), (5,3), (1,4), (2,4), (3,4), (5,4), (6,4), (1,5), (2,5), (3,5), (4,5), (5,5), (1,6), (2,6), (3,6), (4,6), (1,7), (3,7), (4,7), (5,7). Wählen wir davon: P(1,1), O(3,1), T(4,1), S(6,1), D(1,3), A(5,3), M(3,3) (bei mir schon M). Passt das?} \\
            }
        }
    }
    & & \\ \\

    \theutterance \stepcounter{utterance}  
    & & & \multicolumn{2}{p{0.3\linewidth}}{
        \cellcolor[rgb]{0.9,0.9,0.9}{
            \makecell[{{p{\linewidth}}}]{
                \texttt{\tiny{[GM$|$GM]}}
                \texttt{SPRICH: Mein Zeile 2: (1,2) '╟', (2,2) '─', (3,2) '┴', (4,2) '─', (5,2) '┼', (6,2) '─', (7,2) '┤'. Keine leeren Zellen in Zeile 2. Zeile 1: (1,1) '◌', (2,1) '◌', (3,1) '◌', (4,1) '◌', (5,1) '◌', (6,1) '◌', (7,1) '║'. Also leere Zellen nur in Zeile 1: (1,1){-}(6,1). Zeile 3: (1,3) '◌', (2,3) '◌', (3,3) 'M', (4,3) '│', (5,3) '◌', (6,3) '└', (7,3) '─'. Leer: (1,3), (2,3), (5,3). Zeile 4: (1,4) '◌' (war A), (2,4) '◌', (3,4) '◌', (4,4) 'S', (5,4) '◌', (6,4) '◌', (7,4) '║'. Leer: (1,4), (2,4), (3,4), (5,4), (6,4). Zeile 5: (1,5) '◌', (2,5) '◌', (3,5) '◌', (4,5) '◌', (5,5) '◌', (6,5) '┌', (7,5) '─'. Leer: (1,5){-}(5,5). Zeile 6: (1,6) '◌', (2,6) '◌', (3,6) '◌', (4,6) '◌', (5,6) 'P', (6,6) 'T', (7,6) 'O'. Leer: (1,6){-}(4,6). Zeile 7: (1,7) '◌', (2,7) 'D', (3,7) '◌', (4,7) '◌', (5,7) '◌', (6,7) '│', (7,7) '╝'. Leer: (1,7), (3,7), (4,7), (5,7). Gemeinsame leere Zellen mit dir: (1,1), (3,1), (4,1), (6,1), (1,3), (5,3), (1,4), (2,4), (3,4), (5,4), (6,4), (1,5), (2,5), (3,5), (4,5), (5,5), (1,6), (2,6), (3,6), (4,6), (1,7), (3,7), (4,7), (5,7). Wählen wir davon: P(1,1), O(3,1), T(4,1), S(6,1), D(1,3), A(5,3), M(3,3) (bei mir schon M). Passt das?} \\
            }
        }
    }
    & & \\ \\

    \theutterance \stepcounter{utterance}  
    & & \multicolumn{4}{p{0.6\linewidth}}{
        \cellcolor[rgb]{0.9,0.9,0.9}{
            \makecell[{{p{\linewidth}}}]{
                \texttt{\tiny{[GM$\rangle$P2]}}
                \texttt{Deine Nachricht wurde an deinen Mitspieler übermittelt.} \\
\\ 
\\ 
\texttt{Ihr seid gerade bei Runde 10 von maximal 28.} \\
\texttt{Ihr habt gemeinsam 4 von maximal 16 Strafpunkten angesammelt.} \\
\texttt{Dein Mitspieler hat dir diese Nachricht geschickt:} \\
\\ 
\\ 
\texttt{"Mein Zeile 2: (1,2) '╟', (2,2) '─', (3,2) '┴', (4,2) '─', (5,2) '┼', (6,2) '─', (7,2) '┤'. Keine leeren Zellen in Zeile 2. Zeile 1: (1,1) '◌', (2,1) '◌', (3,1) '◌', (4,1) '◌', (5,1) '◌', (6,1) '◌', (7,1) '║'. Also leere Zellen nur in Zeile 1: (1,1){-}(6,1). Zeile 3: (1,3) '◌', (2,3) '◌', (3,3) 'M', (4,3) '│', (5,3) '◌', (6,3) '└', (7,3) '─'. Leer: (1,3), (2,3), (5,3). Zeile 4: (1,4) '◌' (war A), (2,4) '◌', (3,4) '◌', (4,4) 'S', (5,4) '◌', (6,4) '◌', (7,4) '║'. Leer: (1,4), (2,4), (3,4), (5,4), (6,4). Zeile 5: (1,5) '◌', (2,5) '◌', (3,5) '◌', (4,5) '◌', (5,5) '◌', (6,5) '┌', (7,5) '─'. Leer: (1,5){-}(5,5). Zeile 6: (1,6) '◌', (2,6) '◌', (3,6) '◌', (4,6) '◌', (5,6) 'P', (6,6) 'T', (7,6) 'O'. Leer: (1,6){-}(4,6). Zeile 7: (1,7) '◌', (2,7) 'D', (3,7) '◌', (4,7) '◌', (5,7) '◌', (6,7) '│', (7,7) '╝'. Leer: (1,7), (3,7), (4,7), (5,7). Gemeinsame leere Zellen mit dir: (1,1), (3,1), (4,1), (6,1), (1,3), (5,3), (1,4), (2,4), (3,4), (5,4), (6,4), (1,5), (2,5), (3,5), (4,5), (5,5), (1,6), (2,6), (3,6), (4,6), (1,7), (3,7), (4,7), (5,7). Wählen wir davon: P(1,1), O(3,1), T(4,1), S(6,1), D(1,3), A(5,3), M(3,3) (bei mir schon M). Passt das?"} \\
\\ 
\\ 
\texttt{Wie lautet dein nächster Befehl?} \\
            }
        }
    }
    & \\ \\

    \theutterance \stepcounter{utterance}  
    & & & \multicolumn{4}{p{0.6\linewidth}}{
        \cellcolor[rgb]{0.9,0.9,0.9}{
            \makecell[{{p{\linewidth}}}]{
                \texttt{\tiny{[GM$\langle$P2]}}
                \texttt{SPRICH: Perfekt! Das passt: P(1,1), O(3,1), T(4,1), S(6,1), D(1,3), A(5,3), M(3,3). Ich beginne mit Verschiebungen. Zuerst bewege ich mein P von (4,4) nach (1,1).} \\
            }
        }
    }
     \\ \\

    \theutterance \stepcounter{utterance}  
    & & & \multicolumn{2}{p{0.3\linewidth}}{
        \cellcolor[rgb]{0.9,0.9,0.9}{
            \makecell[{{p{\linewidth}}}]{
                \texttt{\tiny{[GM$|$GM]}}
                \texttt{SPRICH: Perfekt! Das passt: P(1,1), O(3,1), T(4,1), S(6,1), D(1,3), A(5,3), M(3,3). Ich beginne mit Verschiebungen. Zuerst bewege ich mein P von (4,4) nach (1,1).} \\
            }
        }
    }
    & & \\ \\

    \theutterance \stepcounter{utterance}  
    & & \multicolumn{4}{p{0.6\linewidth}}{
        \cellcolor[rgb]{0.9,0.9,0.9}{
            \makecell[{{p{\linewidth}}}]{
                \texttt{\tiny{[P1$\langle$GM]}}
                \texttt{Deine Nachricht wurde an deinen Mitspieler übermittelt.} \\
\\ 
\\ 
\texttt{Ihr seid gerade bei Runde 10 von maximal 28.} \\
\texttt{Ihr habt gemeinsam 4 von maximal 16 Strafpunkten angesammelt.} \\
\texttt{Dein Mitspieler hat dir diese Nachricht geschickt:} \\
\\ 
\\ 
\texttt{"Perfekt! Das passt: P(1,1), O(3,1), T(4,1), S(6,1), D(1,3), A(5,3), M(3,3). Ich beginne mit Verschiebungen. Zuerst bewege ich mein P von (4,4) nach (1,1)."} \\
\\ 
\\ 
\texttt{Wie lautet dein nächster Befehl?} \\
            }
        }
    }
    & \\ \\

    \theutterance \stepcounter{utterance}  
    & \multicolumn{4}{p{0.6\linewidth}}{
        \cellcolor[rgb]{0.9,0.9,0.9}{
            \makecell[{{p{\linewidth}}}]{
                \texttt{\tiny{[P1$\rangle$GM]}}
                \texttt{VERSCHIEBE: A, (5,3)} \\
            }
        }
    }
    & & \\ \\

    \theutterance \stepcounter{utterance}  
    & & & \multicolumn{2}{p{0.3\linewidth}}{
        \cellcolor[rgb]{0.9,0.9,0.9}{
            \makecell[{{p{\linewidth}}}]{
                \texttt{\tiny{[GM$|$GM]}}
                \texttt{VERSCHIEBE: A, (5,3)} \\
            }
        }
    }
    & & \\ \\

    \theutterance \stepcounter{utterance}  
    & & & \multicolumn{2}{p{0.3\linewidth}}{
        \cellcolor[rgb]{0.9,0.9,0.9}{
            \makecell[{{p{\linewidth}}}]{
                \texttt{\tiny{[GM$|$GM]}}
                \texttt{'A' erfolgreich nach (5,3) verschoben. Dein Raster sieht jetzt so aus:} \\
\texttt{```} \\
\texttt{ 1234567} \\
\texttt{╔═╤═╤═╤═╗} \\
\texttt{║◌│◌│◌│◌║ 1} \\
\texttt{╟─┴─┼─┤◌║ 2} \\
\texttt{║◌◌M│A└─╢ 3} \\
\texttt{╟───┘◌S◌║ 4} \\
\texttt{║◌◌◌◌◌┌─╢ 5} \\
\texttt{║◌◌◌PT│O║ 6} \\
\texttt{║◌D◌◌◌│◌║ 7} \\
\texttt{╚═════╧═╝} \\
\texttt{```} \\
            }
        }
    }
    & & \\ \\

    \theutterance \stepcounter{utterance}  
    & & \multicolumn{4}{p{0.6\linewidth}}{
        \cellcolor[rgb]{0.9,0.9,0.9}{
            \makecell[{{p{\linewidth}}}]{
                \texttt{\tiny{[GM$\rangle$P2]}}
                \texttt{Deine Nachricht wurde an deinen Mitspieler übermittelt.} \\
\\ 
\\ 
\texttt{Ihr seid gerade bei Runde 11 von maximal 28.} \\
\texttt{Ihr habt gemeinsam 4 von maximal 16 Strafpunkten angesammelt.} \\
\texttt{Dein Mitspieler hat ein Objekt auf seinem Raster bewegt.} \\
\\ 
\\ 
\texttt{Wie lautet dein nächster Befehl?} \\
            }
        }
    }
    & \\ \\

    \theutterance \stepcounter{utterance}  
    & & & \multicolumn{4}{p{0.6\linewidth}}{
        \cellcolor[rgb]{0.9,0.9,0.9}{
            \makecell[{{p{\linewidth}}}]{
                \texttt{\tiny{[GM$\langle$P2]}}
                \texttt{VERSCHIEBE: P, (1,1)} \\
            }
        }
    }
     \\ \\

    \theutterance \stepcounter{utterance}  
    & & & \multicolumn{2}{p{0.3\linewidth}}{
        \cellcolor[rgb]{0.9,0.9,0.9}{
            \makecell[{{p{\linewidth}}}]{
                \texttt{\tiny{[GM$|$GM]}}
                \texttt{VERSCHIEBE: P, (1,1)} \\
            }
        }
    }
    & & \\ \\

    \theutterance \stepcounter{utterance}  
    & & & \multicolumn{2}{p{0.3\linewidth}}{
        \cellcolor[rgb]{0.9,0.9,0.9}{
            \makecell[{{p{\linewidth}}}]{
                \texttt{\tiny{[GM$|$GM]}}
                \texttt{'P' erfolgreich nach (1,1) verschoben. Dein Raster sieht jetzt so aus:} \\
\texttt{```} \\
\texttt{ 1234567} \\
\texttt{╔═╤═╤═╤═╗} \\
\texttt{║P│◌│◌│◌║ 1} \\
\texttt{╟─┴─┼─┤◌║ 2} \\
\texttt{║◌M◌│◌└─╢ 3} \\
\texttt{╟───┘◌◌◌║ 4} \\
\texttt{║O◌◌◌◌┌─╢ 5} \\
\texttt{║◌D◌◌A│S║ 6} \\
\texttt{║◌◌◌◌T│◌║ 7} \\
\texttt{╚═════╧═╝} \\
\texttt{```} \\
            }
        }
    }
    & & \\ \\

    \theutterance \stepcounter{utterance}  
    & & \multicolumn{4}{p{0.6\linewidth}}{
        \cellcolor[rgb]{0.9,0.9,0.9}{
            \makecell[{{p{\linewidth}}}]{
                \texttt{\tiny{[P1$\langle$GM]}}
                \texttt{'A' erfolgreich nach (5,3) verschoben. Dein Raster sieht jetzt so aus:} \\
\texttt{```} \\
\texttt{ 1234567} \\
\texttt{╔═╤═╤═╤═╗} \\
\texttt{║◌│◌│◌│◌║ 1} \\
\texttt{╟─┴─┼─┤◌║ 2} \\
\texttt{║◌◌M│A└─╢ 3} \\
\texttt{╟───┘◌S◌║ 4} \\
\texttt{║◌◌◌◌◌┌─╢ 5} \\
\texttt{║◌◌◌PT│O║ 6} \\
\texttt{║◌D◌◌◌│◌║ 7} \\
\texttt{╚═════╧═╝} \\
\texttt{```} \\
\texttt{Ihr seid gerade bei Runde 11 von maximal 28.} \\
\texttt{Ihr habt gemeinsam 4 von maximal 16 Strafpunkten angesammelt.} \\
\texttt{Dein Mitspieler hat ein Objekt auf seinem Raster bewegt.} \\
\\ 
\\ 
\texttt{Wie lautet dein nächster Befehl?} \\
            }
        }
    }
    & \\ \\

    \theutterance \stepcounter{utterance}  
    & \multicolumn{4}{p{0.6\linewidth}}{
        \cellcolor[rgb]{0.9,0.9,0.9}{
            \makecell[{{p{\linewidth}}}]{
                \texttt{\tiny{[P1$\rangle$GM]}}
                \texttt{VERSCHIEBE: D, (1,3)} \\
            }
        }
    }
    & & \\ \\

    \theutterance \stepcounter{utterance}  
    & & & \multicolumn{2}{p{0.3\linewidth}}{
        \cellcolor[rgb]{0.9,0.9,0.9}{
            \makecell[{{p{\linewidth}}}]{
                \texttt{\tiny{[GM$|$GM]}}
                \texttt{VERSCHIEBE: D, (1,3)} \\
            }
        }
    }
    & & \\ \\

    \theutterance \stepcounter{utterance}  
    & & & \multicolumn{2}{p{0.3\linewidth}}{
        \cellcolor[rgb]{0.9,0.9,0.9}{
            \makecell[{{p{\linewidth}}}]{
                \texttt{\tiny{[GM$|$GM]}}
                \texttt{'D' erfolgreich nach (1,3) verschoben. Dein Raster sieht jetzt so aus:} \\
\texttt{```} \\
\texttt{ 1234567} \\
\texttt{╔═╤═╤═╤═╗} \\
\texttt{║◌│◌│◌│◌║ 1} \\
\texttt{╟─┴─┼─┤◌║ 2} \\
\texttt{║D◌M│A└─╢ 3} \\
\texttt{╟───┘◌S◌║ 4} \\
\texttt{║◌◌◌◌◌┌─╢ 5} \\
\texttt{║◌◌◌PT│O║ 6} \\
\texttt{║◌◌◌◌◌│◌║ 7} \\
\texttt{╚═════╧═╝} \\
\texttt{```} \\
            }
        }
    }
    & & \\ \\

    \theutterance \stepcounter{utterance}  
    & & \multicolumn{4}{p{0.6\linewidth}}{
        \cellcolor[rgb]{0.9,0.9,0.9}{
            \makecell[{{p{\linewidth}}}]{
                \texttt{\tiny{[GM$\rangle$P2]}}
                \texttt{'P' erfolgreich nach (1,1) verschoben. Dein Raster sieht jetzt so aus:} \\
\texttt{```} \\
\texttt{ 1234567} \\
\texttt{╔═╤═╤═╤═╗} \\
\texttt{║P│◌│◌│◌║ 1} \\
\texttt{╟─┴─┼─┤◌║ 2} \\
\texttt{║◌M◌│◌└─╢ 3} \\
\texttt{╟───┘◌◌◌║ 4} \\
\texttt{║O◌◌◌◌┌─╢ 5} \\
\texttt{║◌D◌◌A│S║ 6} \\
\texttt{║◌◌◌◌T│◌║ 7} \\
\texttt{╚═════╧═╝} \\
\texttt{```} \\
\texttt{Ihr seid gerade bei Runde 12 von maximal 28.} \\
\texttt{Ihr habt gemeinsam 4 von maximal 16 Strafpunkten angesammelt.} \\
\texttt{Dein Mitspieler hat ein Objekt auf seinem Raster bewegt.} \\
\\ 
\\ 
\texttt{Wie lautet dein nächster Befehl?} \\
            }
        }
    }
    & \\ \\

    \theutterance \stepcounter{utterance}  
    & & & \multicolumn{4}{p{0.6\linewidth}}{
        \cellcolor[rgb]{0.9,0.9,0.9}{
            \makecell[{{p{\linewidth}}}]{
                \texttt{\tiny{[GM$\langle$P2]}}
                \texttt{VERSCHIEBE: O, (3,1)} \\
            }
        }
    }
     \\ \\

    \theutterance \stepcounter{utterance}  
    & & & \multicolumn{2}{p{0.3\linewidth}}{
        \cellcolor[rgb]{0.9,0.9,0.9}{
            \makecell[{{p{\linewidth}}}]{
                \texttt{\tiny{[GM$|$GM]}}
                \texttt{VERSCHIEBE: O, (3,1)} \\
            }
        }
    }
    & & \\ \\

    \theutterance \stepcounter{utterance}  
    & & & \multicolumn{2}{p{0.3\linewidth}}{
        \cellcolor[rgb]{0.9,0.9,0.9}{
            \makecell[{{p{\linewidth}}}]{
                \texttt{\tiny{[GM$|$GM]}}
                \texttt{'O' erfolgreich nach (3,1) verschoben. Dein Raster sieht jetzt so aus:} \\
\texttt{```} \\
\texttt{ 1234567} \\
\texttt{╔═╤═╤═╤═╗} \\
\texttt{║P│O│◌│◌║ 1} \\
\texttt{╟─┴─┼─┤◌║ 2} \\
\texttt{║◌M◌│◌└─╢ 3} \\
\texttt{╟───┘◌◌◌║ 4} \\
\texttt{║◌◌◌◌◌┌─╢ 5} \\
\texttt{║◌D◌◌A│S║ 6} \\
\texttt{║◌◌◌◌T│◌║ 7} \\
\texttt{╚═════╧═╝} \\
\texttt{```} \\
            }
        }
    }
    & & \\ \\

    \theutterance \stepcounter{utterance}  
    & & \multicolumn{4}{p{0.6\linewidth}}{
        \cellcolor[rgb]{0.9,0.9,0.9}{
            \makecell[{{p{\linewidth}}}]{
                \texttt{\tiny{[P1$\langle$GM]}}
                \texttt{'D' erfolgreich nach (1,3) verschoben. Dein Raster sieht jetzt so aus:} \\
\texttt{```} \\
\texttt{ 1234567} \\
\texttt{╔═╤═╤═╤═╗} \\
\texttt{║◌│◌│◌│◌║ 1} \\
\texttt{╟─┴─┼─┤◌║ 2} \\
\texttt{║D◌M│A└─╢ 3} \\
\texttt{╟───┘◌S◌║ 4} \\
\texttt{║◌◌◌◌◌┌─╢ 5} \\
\texttt{║◌◌◌PT│O║ 6} \\
\texttt{║◌◌◌◌◌│◌║ 7} \\
\texttt{╚═════╧═╝} \\
\texttt{```} \\
\texttt{Ihr seid gerade bei Runde 12 von maximal 28.} \\
\texttt{Ihr habt gemeinsam 4 von maximal 16 Strafpunkten angesammelt.} \\
\texttt{Dein Mitspieler hat ein Objekt auf seinem Raster bewegt.} \\
\\ 
\\ 
\texttt{Wie lautet dein nächster Befehl?} \\
            }
        }
    }
    & \\ \\

    \theutterance \stepcounter{utterance}  
    & \multicolumn{4}{p{0.6\linewidth}}{
        \cellcolor[rgb]{0.9,0.9,0.9}{
            \makecell[{{p{\linewidth}}}]{
                \texttt{\tiny{[P1$\rangle$GM]}}
                \texttt{VERSCHIEBE: P, (1,1)} \\
            }
        }
    }
    & & \\ \\

    \theutterance \stepcounter{utterance}  
    & & & \multicolumn{2}{p{0.3\linewidth}}{
        \cellcolor[rgb]{0.9,0.9,0.9}{
            \makecell[{{p{\linewidth}}}]{
                \texttt{\tiny{[GM$|$GM]}}
                \texttt{VERSCHIEBE: P, (1,1)} \\
            }
        }
    }
    & & \\ \\

    \theutterance \stepcounter{utterance}  
    & & & \multicolumn{2}{p{0.3\linewidth}}{
        \cellcolor[rgb]{0.9,0.9,0.9}{
            \makecell[{{p{\linewidth}}}]{
                \texttt{\tiny{[GM$|$GM]}}
                \texttt{'P' erfolgreich nach (1,1) verschoben. Dein Raster sieht jetzt so aus:} \\
\texttt{```} \\
\texttt{ 1234567} \\
\texttt{╔═╤═╤═╤═╗} \\
\texttt{║P│◌│◌│◌║ 1} \\
\texttt{╟─┴─┼─┤◌║ 2} \\
\texttt{║D◌M│A└─╢ 3} \\
\texttt{╟───┘◌S◌║ 4} \\
\texttt{║◌◌◌◌◌┌─╢ 5} \\
\texttt{║◌◌◌◌T│O║ 6} \\
\texttt{║◌◌◌◌◌│◌║ 7} \\
\texttt{╚═════╧═╝} \\
\texttt{```} \\
            }
        }
    }
    & & \\ \\

    \theutterance \stepcounter{utterance}  
    & & \multicolumn{4}{p{0.6\linewidth}}{
        \cellcolor[rgb]{0.9,0.9,0.9}{
            \makecell[{{p{\linewidth}}}]{
                \texttt{\tiny{[GM$\rangle$P2]}}
                \texttt{'O' erfolgreich nach (3,1) verschoben. Dein Raster sieht jetzt so aus:} \\
\texttt{```} \\
\texttt{ 1234567} \\
\texttt{╔═╤═╤═╤═╗} \\
\texttt{║P│O│◌│◌║ 1} \\
\texttt{╟─┴─┼─┤◌║ 2} \\
\texttt{║◌M◌│◌└─╢ 3} \\
\texttt{╟───┘◌◌◌║ 4} \\
\texttt{║◌◌◌◌◌┌─╢ 5} \\
\texttt{║◌D◌◌A│S║ 6} \\
\texttt{║◌◌◌◌T│◌║ 7} \\
\texttt{╚═════╧═╝} \\
\texttt{```} \\
\texttt{Ihr seid gerade bei Runde 13 von maximal 28.} \\
\texttt{Ihr habt gemeinsam 4 von maximal 16 Strafpunkten angesammelt.} \\
\texttt{Dein Mitspieler hat ein Objekt auf seinem Raster bewegt.} \\
\\ 
\\ 
\texttt{Wie lautet dein nächster Befehl?} \\
            }
        }
    }
    & \\ \\

    \theutterance \stepcounter{utterance}  
    & & & \multicolumn{4}{p{0.6\linewidth}}{
        \cellcolor[rgb]{0.9,0.9,0.9}{
            \makecell[{{p{\linewidth}}}]{
                \texttt{\tiny{[GM$\langle$P2]}}
                \texttt{VERSCHIEBE: T, (4,1)} \\
            }
        }
    }
     \\ \\

    \theutterance \stepcounter{utterance}  
    & & & \multicolumn{2}{p{0.3\linewidth}}{
        \cellcolor[rgb]{0.9,0.9,0.9}{
            \makecell[{{p{\linewidth}}}]{
                \texttt{\tiny{[GM$|$GM]}}
                \texttt{VERSCHIEBE: T, (4,1)} \\
            }
        }
    }
    & & \\ \\

    \theutterance \stepcounter{utterance}  
    & & & \multicolumn{2}{p{0.3\linewidth}}{
        \cellcolor[rgb]{0.9,0.9,0.9}{
            \makecell[{{p{\linewidth}}}]{
                \texttt{\tiny{[GM$|$GM]}}
                \texttt{Strafpunkt: (4,1) ist nicht leer, sondern enthält '│'.} \\
\\ 
\texttt{Ihr habt gemeinsam 5 von maximal 16 Strafpunkten angesammelt.} \\
\texttt{Bitte versuche es noch mal!} \\
            }
        }
    }
    & & \\ \\

    \theutterance \stepcounter{utterance}  
    & & \multicolumn{4}{p{0.6\linewidth}}{
        \cellcolor[rgb]{0.9,0.9,0.9}{
            \makecell[{{p{\linewidth}}}]{
                \texttt{\tiny{[GM$\rangle$P2]}}
                \texttt{Strafpunkt: (4,1) ist nicht leer, sondern enthält '│'.} \\
\\ 
\texttt{Ihr habt gemeinsam 5 von maximal 16 Strafpunkten angesammelt.} \\
\texttt{Bitte versuche es noch mal!} \\
            }
        }
    }
    & \\ \\

    \theutterance \stepcounter{utterance}  
    & & & \multicolumn{4}{p{0.6\linewidth}}{
        \cellcolor[rgb]{0.9,0.9,0.9}{
            \makecell[{{p{\linewidth}}}]{
                \texttt{\tiny{[GM$\langle$P2]}}
                \texttt{SPRICH: Bei mir ist (4,1) '│' (Rahmenteil). Wir müssen T an eine andere leere Stelle setzen. Vorschlag: T(6,1) stattdessen, und S auf (4,1) überspringen? Aber (4,1) ist nicht leer. Korrigiere Ziel für T: T(6,1) und S(4,5) oder ähnlich? Oder wir ändern den Plan: Da (4,1) bei mir nicht leer ist, nutzen wir T(5,1) – ist das bei dir leer? Bei mir ist (5,1) jetzt leer (nachdem O verschoben wurde). Also neue Ziele: P(1,1), O(3,1), T(5,1), S(6,1), D(1,3), A(5,3), M(3,3). Bestätige.} \\
            }
        }
    }
     \\ \\

    \theutterance \stepcounter{utterance}  
    & & & \multicolumn{2}{p{0.3\linewidth}}{
        \cellcolor[rgb]{0.9,0.9,0.9}{
            \makecell[{{p{\linewidth}}}]{
                \texttt{\tiny{[GM$|$GM]}}
                \texttt{SPRICH: Bei mir ist (4,1) '│' (Rahmenteil). Wir müssen T an eine andere leere Stelle setzen. Vorschlag: T(6,1) stattdessen, und S auf (4,1) überspringen? Aber (4,1) ist nicht leer. Korrigiere Ziel für T: T(6,1) und S(4,5) oder ähnlich? Oder wir ändern den Plan: Da (4,1) bei mir nicht leer ist, nutzen wir T(5,1) – ist das bei dir leer? Bei mir ist (5,1) jetzt leer (nachdem O verschoben wurde). Also neue Ziele: P(1,1), O(3,1), T(5,1), S(6,1), D(1,3), A(5,3), M(3,3). Bestätige.} \\
            }
        }
    }
    & & \\ \\

    \theutterance \stepcounter{utterance}  
    & & \multicolumn{4}{p{0.6\linewidth}}{
        \cellcolor[rgb]{0.9,0.9,0.9}{
            \makecell[{{p{\linewidth}}}]{
                \texttt{\tiny{[P1$\langle$GM]}}
                \texttt{'P' erfolgreich nach (1,1) verschoben. Dein Raster sieht jetzt so aus:} \\
\texttt{```} \\
\texttt{ 1234567} \\
\texttt{╔═╤═╤═╤═╗} \\
\texttt{║P│◌│◌│◌║ 1} \\
\texttt{╟─┴─┼─┤◌║ 2} \\
\texttt{║D◌M│A└─╢ 3} \\
\texttt{╟───┘◌S◌║ 4} \\
\texttt{║◌◌◌◌◌┌─╢ 5} \\
\texttt{║◌◌◌◌T│O║ 6} \\
\texttt{║◌◌◌◌◌│◌║ 7} \\
\texttt{╚═════╧═╝} \\
\texttt{```} \\
\texttt{Ihr seid gerade bei Runde 13 von maximal 28.} \\
\texttt{Ihr habt gemeinsam 5 von maximal 16 Strafpunkten angesammelt.} \\
\texttt{Dein Mitspieler hat dir diese Nachricht geschickt:} \\
\\ 
\\ 
\texttt{"Bei mir ist (4,1) '│' (Rahmenteil). Wir müssen T an eine andere leere Stelle setzen. Vorschlag: T(6,1) stattdessen, und S auf (4,1) überspringen? Aber (4,1) ist nicht leer. Korrigiere Ziel für T: T(6,1) und S(4,5) oder ähnlich? Oder wir ändern den Plan: Da (4,1) bei mir nicht leer ist, nutzen wir T(5,1) – ist das bei dir leer? Bei mir ist (5,1) jetzt leer (nachdem O verschoben wurde). Also neue Ziele: P(1,1), O(3,1), T(5,1), S(6,1), D(1,3), A(5,3), M(3,3). Bestätige."} \\
\\ 
\\ 
\texttt{Wie lautet dein nächster Befehl?} \\
            }
        }
    }
    & \\ \\

    \theutterance \stepcounter{utterance}  
    & \multicolumn{4}{p{0.6\linewidth}}{
        \cellcolor[rgb]{0.9,0.9,0.9}{
            \makecell[{{p{\linewidth}}}]{
                \texttt{\tiny{[P1$\rangle$GM]}}
                \texttt{SPRICH: Bestätige neue Ziele: P(1,1), O(3,1), T(5,1), S(6,1), D(1,3), A(5,3), M(3,3). (5,1) ist bei mir leer. Ich verschiebe O nach (3,1).} \\
            }
        }
    }
    & & \\ \\

    \theutterance \stepcounter{utterance}  
    & & & \multicolumn{2}{p{0.3\linewidth}}{
        \cellcolor[rgb]{0.9,0.9,0.9}{
            \makecell[{{p{\linewidth}}}]{
                \texttt{\tiny{[GM$|$GM]}}
                \texttt{SPRICH: Bestätige neue Ziele: P(1,1), O(3,1), T(5,1), S(6,1), D(1,3), A(5,3), M(3,3). (5,1) ist bei mir leer. Ich verschiebe O nach (3,1).} \\
            }
        }
    }
    & & \\ \\

    \theutterance \stepcounter{utterance}  
    & & \multicolumn{4}{p{0.6\linewidth}}{
        \cellcolor[rgb]{0.9,0.9,0.9}{
            \makecell[{{p{\linewidth}}}]{
                \texttt{\tiny{[GM$\rangle$P2]}}
                \texttt{Deine Nachricht wurde an deinen Mitspieler übermittelt.} \\
\\ 
\\ 
\texttt{Ihr seid gerade bei Runde 14 von maximal 28.} \\
\texttt{Ihr habt gemeinsam 5 von maximal 16 Strafpunkten angesammelt.} \\
\texttt{Dein Mitspieler hat dir diese Nachricht geschickt:} \\
\\ 
\\ 
\texttt{"Bestätige neue Ziele: P(1,1), O(3,1), T(5,1), S(6,1), D(1,3), A(5,3), M(3,3). (5,1) ist bei mir leer. Ich verschiebe O nach (3,1)."} \\
\\ 
\\ 
\texttt{Wie lautet dein nächster Befehl?} \\
            }
        }
    }
    & \\ \\

    \theutterance \stepcounter{utterance}  
    & & & \multicolumn{4}{p{0.6\linewidth}}{
        \cellcolor[rgb]{0.9,0.9,0.9}{
            \makecell[{{p{\linewidth}}}]{
                \texttt{\tiny{[GM$\langle$P2]}}
                \texttt{VERSCHIEBE: T, (5,1)} \\
            }
        }
    }
     \\ \\

    \theutterance \stepcounter{utterance}  
    & & & \multicolumn{2}{p{0.3\linewidth}}{
        \cellcolor[rgb]{0.9,0.9,0.9}{
            \makecell[{{p{\linewidth}}}]{
                \texttt{\tiny{[GM$|$GM]}}
                \texttt{VERSCHIEBE: T, (5,1)} \\
            }
        }
    }
    & & \\ \\

    \theutterance \stepcounter{utterance}  
    & & & \multicolumn{2}{p{0.3\linewidth}}{
        \cellcolor[rgb]{0.9,0.9,0.9}{
            \makecell[{{p{\linewidth}}}]{
                \texttt{\tiny{[GM$|$GM]}}
                \texttt{'T' erfolgreich nach (5,1) verschoben. Dein Raster sieht jetzt so aus:} \\
\texttt{```} \\
\texttt{ 1234567} \\
\texttt{╔═╤═╤═╤═╗} \\
\texttt{║P│O│T│◌║ 1} \\
\texttt{╟─┴─┼─┤◌║ 2} \\
\texttt{║◌M◌│◌└─╢ 3} \\
\texttt{╟───┘◌◌◌║ 4} \\
\texttt{║◌◌◌◌◌┌─╢ 5} \\
\texttt{║◌D◌◌A│S║ 6} \\
\texttt{║◌◌◌◌◌│◌║ 7} \\
\texttt{╚═════╧═╝} \\
\texttt{```} \\
            }
        }
    }
    & & \\ \\

    \theutterance \stepcounter{utterance}  
    & & \multicolumn{4}{p{0.6\linewidth}}{
        \cellcolor[rgb]{0.9,0.9,0.9}{
            \makecell[{{p{\linewidth}}}]{
                \texttt{\tiny{[P1$\langle$GM]}}
                \texttt{Deine Nachricht wurde an deinen Mitspieler übermittelt.} \\
\\ 
\\ 
\texttt{Ihr seid gerade bei Runde 14 von maximal 28.} \\
\texttt{Ihr habt gemeinsam 5 von maximal 16 Strafpunkten angesammelt.} \\
\texttt{Dein Mitspieler hat ein Objekt auf seinem Raster bewegt.} \\
\\ 
\\ 
\texttt{Wie lautet dein nächster Befehl?} \\
            }
        }
    }
    & \\ \\

    \theutterance \stepcounter{utterance}  
    & \multicolumn{4}{p{0.6\linewidth}}{
        \cellcolor[rgb]{0.9,0.9,0.9}{
            \makecell[{{p{\linewidth}}}]{
                \texttt{\tiny{[P1$\rangle$GM]}}
                \texttt{VERSCHIEBE: O, (3,1)} \\
            }
        }
    }
    & & \\ \\

    \theutterance \stepcounter{utterance}  
    & & & \multicolumn{2}{p{0.3\linewidth}}{
        \cellcolor[rgb]{0.9,0.9,0.9}{
            \makecell[{{p{\linewidth}}}]{
                \texttt{\tiny{[GM$|$GM]}}
                \texttt{VERSCHIEBE: O, (3,1)} \\
            }
        }
    }
    & & \\ \\

    \theutterance \stepcounter{utterance}  
    & & & \multicolumn{2}{p{0.3\linewidth}}{
        \cellcolor[rgb]{0.9,0.9,0.9}{
            \makecell[{{p{\linewidth}}}]{
                \texttt{\tiny{[GM$|$GM]}}
                \texttt{'O' erfolgreich nach (3,1) verschoben. Dein Raster sieht jetzt so aus:} \\
\texttt{```} \\
\texttt{ 1234567} \\
\texttt{╔═╤═╤═╤═╗} \\
\texttt{║P│O│◌│◌║ 1} \\
\texttt{╟─┴─┼─┤◌║ 2} \\
\texttt{║D◌M│A└─╢ 3} \\
\texttt{╟───┘◌S◌║ 4} \\
\texttt{║◌◌◌◌◌┌─╢ 5} \\
\texttt{║◌◌◌◌T│◌║ 6} \\
\texttt{║◌◌◌◌◌│◌║ 7} \\
\texttt{╚═════╧═╝} \\
\texttt{```} \\
            }
        }
    }
    & & \\ \\

    \theutterance \stepcounter{utterance}  
    & & \multicolumn{4}{p{0.6\linewidth}}{
        \cellcolor[rgb]{0.9,0.9,0.9}{
            \makecell[{{p{\linewidth}}}]{
                \texttt{\tiny{[GM$\rangle$P2]}}
                \texttt{'T' erfolgreich nach (5,1) verschoben. Dein Raster sieht jetzt so aus:} \\
\texttt{```} \\
\texttt{ 1234567} \\
\texttt{╔═╤═╤═╤═╗} \\
\texttt{║P│O│T│◌║ 1} \\
\texttt{╟─┴─┼─┤◌║ 2} \\
\texttt{║◌M◌│◌└─╢ 3} \\
\texttt{╟───┘◌◌◌║ 4} \\
\texttt{║◌◌◌◌◌┌─╢ 5} \\
\texttt{║◌D◌◌A│S║ 6} \\
\texttt{║◌◌◌◌◌│◌║ 7} \\
\texttt{╚═════╧═╝} \\
\texttt{```} \\
\texttt{Ihr seid gerade bei Runde 15 von maximal 28.} \\
\texttt{Ihr habt gemeinsam 5 von maximal 16 Strafpunkten angesammelt.} \\
\texttt{Dein Mitspieler hat ein Objekt auf seinem Raster bewegt.} \\
\\ 
\\ 
\texttt{Wie lautet dein nächster Befehl?} \\
            }
        }
    }
    & \\ \\

    \theutterance \stepcounter{utterance}  
    & & & \multicolumn{4}{p{0.6\linewidth}}{
        \cellcolor[rgb]{0.9,0.9,0.9}{
            \makecell[{{p{\linewidth}}}]{
                \texttt{\tiny{[GM$\langle$P2]}}
                \texttt{VERSCHIEBE: S, (6,1)} \\
            }
        }
    }
     \\ \\

    \theutterance \stepcounter{utterance}  
    & & & \multicolumn{2}{p{0.3\linewidth}}{
        \cellcolor[rgb]{0.9,0.9,0.9}{
            \makecell[{{p{\linewidth}}}]{
                \texttt{\tiny{[GM$|$GM]}}
                \texttt{VERSCHIEBE: S, (6,1)} \\
            }
        }
    }
    & & \\ \\

    \theutterance \stepcounter{utterance}  
    & & & \multicolumn{2}{p{0.3\linewidth}}{
        \cellcolor[rgb]{0.9,0.9,0.9}{
            \makecell[{{p{\linewidth}}}]{
                \texttt{\tiny{[GM$|$GM]}}
                \texttt{Strafpunkt: (6,1) ist nicht leer, sondern enthält '│'.} \\
\\ 
\texttt{Ihr habt gemeinsam 6 von maximal 16 Strafpunkten angesammelt.} \\
\texttt{Bitte versuche es noch mal!} \\
            }
        }
    }
    & & \\ \\

    \theutterance \stepcounter{utterance}  
    & & \multicolumn{4}{p{0.6\linewidth}}{
        \cellcolor[rgb]{0.9,0.9,0.9}{
            \makecell[{{p{\linewidth}}}]{
                \texttt{\tiny{[GM$\rangle$P2]}}
                \texttt{Strafpunkt: (6,1) ist nicht leer, sondern enthält '│'.} \\
\\ 
\texttt{Ihr habt gemeinsam 6 von maximal 16 Strafpunkten angesammelt.} \\
\texttt{Bitte versuche es noch mal!} \\
            }
        }
    }
    & \\ \\

    \theutterance \stepcounter{utterance}  
    & & & \multicolumn{4}{p{0.6\linewidth}}{
        \cellcolor[rgb]{0.9,0.9,0.9}{
            \makecell[{{p{\linewidth}}}]{
                \texttt{\tiny{[GM$\langle$P2]}}
                \texttt{SPRICH: (6,1) ist bei mir '│' (Rahmenteil). Wir müssen S an eine andere leere Stelle setzen. Vorschlag: S(4,5) – ist das bei dir leer? Bei mir ist (4,5) leer. Also neue Ziele: S(4,5), behalte den Rest: P(1,1), O(3,1), T(5,1), D(1,3), A(5,3), M(3,3). Bestätige.} \\
            }
        }
    }
     \\ \\

    \theutterance \stepcounter{utterance}  
    & & & \multicolumn{2}{p{0.3\linewidth}}{
        \cellcolor[rgb]{0.9,0.9,0.9}{
            \makecell[{{p{\linewidth}}}]{
                \texttt{\tiny{[GM$|$GM]}}
                \texttt{SPRICH: (6,1) ist bei mir '│' (Rahmenteil). Wir müssen S an eine andere leere Stelle setzen. Vorschlag: S(4,5) – ist das bei dir leer? Bei mir ist (4,5) leer. Also neue Ziele: S(4,5), behalte den Rest: P(1,1), O(3,1), T(5,1), D(1,3), A(5,3), M(3,3). Bestätige.} \\
            }
        }
    }
    & & \\ \\

    \theutterance \stepcounter{utterance}  
    & & \multicolumn{4}{p{0.6\linewidth}}{
        \cellcolor[rgb]{0.9,0.9,0.9}{
            \makecell[{{p{\linewidth}}}]{
                \texttt{\tiny{[P1$\langle$GM]}}
                \texttt{'O' erfolgreich nach (3,1) verschoben. Dein Raster sieht jetzt so aus:} \\
\texttt{```} \\
\texttt{ 1234567} \\
\texttt{╔═╤═╤═╤═╗} \\
\texttt{║P│O│◌│◌║ 1} \\
\texttt{╟─┴─┼─┤◌║ 2} \\
\texttt{║D◌M│A└─╢ 3} \\
\texttt{╟───┘◌S◌║ 4} \\
\texttt{║◌◌◌◌◌┌─╢ 5} \\
\texttt{║◌◌◌◌T│◌║ 6} \\
\texttt{║◌◌◌◌◌│◌║ 7} \\
\texttt{╚═════╧═╝} \\
\texttt{```} \\
\texttt{Ihr seid gerade bei Runde 15 von maximal 28.} \\
\texttt{Ihr habt gemeinsam 6 von maximal 16 Strafpunkten angesammelt.} \\
\texttt{Dein Mitspieler hat dir diese Nachricht geschickt:} \\
\\ 
\\ 
\texttt{"(6,1) ist bei mir '│' (Rahmenteil). Wir müssen S an eine andere leere Stelle setzen. Vorschlag: S(4,5) – ist das bei dir leer? Bei mir ist (4,5) leer. Also neue Ziele: S(4,5), behalte den Rest: P(1,1), O(3,1), T(5,1), D(1,3), A(5,3), M(3,3). Bestätige."} \\
\\ 
\\ 
\texttt{Wie lautet dein nächster Befehl?} \\
            }
        }
    }
    & \\ \\

    \theutterance \stepcounter{utterance}  
    & \multicolumn{4}{p{0.6\linewidth}}{
        \cellcolor[rgb]{0.9,0.9,0.9}{
            \makecell[{{p{\linewidth}}}]{
                \texttt{\tiny{[P1$\rangle$GM]}}
                \texttt{SPRICH: Bestätige S(4,5). (4,5) ist bei mir leer. Ich verschiebe T nach (5,1).} \\
            }
        }
    }
    & & \\ \\

    \theutterance \stepcounter{utterance}  
    & & & \multicolumn{2}{p{0.3\linewidth}}{
        \cellcolor[rgb]{0.9,0.9,0.9}{
            \makecell[{{p{\linewidth}}}]{
                \texttt{\tiny{[GM$|$GM]}}
                \texttt{SPRICH: Bestätige S(4,5). (4,5) ist bei mir leer. Ich verschiebe T nach (5,1).} \\
            }
        }
    }
    & & \\ \\

    \theutterance \stepcounter{utterance}  
    & & \multicolumn{4}{p{0.6\linewidth}}{
        \cellcolor[rgb]{0.9,0.9,0.9}{
            \makecell[{{p{\linewidth}}}]{
                \texttt{\tiny{[GM$\rangle$P2]}}
                \texttt{Deine Nachricht wurde an deinen Mitspieler übermittelt.} \\
\\ 
\\ 
\texttt{Ihr seid gerade bei Runde 16 von maximal 28.} \\
\texttt{Ihr habt gemeinsam 6 von maximal 16 Strafpunkten angesammelt.} \\
\texttt{Dein Mitspieler hat dir diese Nachricht geschickt:} \\
\\ 
\\ 
\texttt{"Bestätige S(4,5). (4,5) ist bei mir leer. Ich verschiebe T nach (5,1)."} \\
\\ 
\\ 
\texttt{Wie lautet dein nächster Befehl?} \\
            }
        }
    }
    & \\ \\

    \theutterance \stepcounter{utterance}  
    & & & \multicolumn{4}{p{0.6\linewidth}}{
        \cellcolor[rgb]{0.9,0.9,0.9}{
            \makecell[{{p{\linewidth}}}]{
                \texttt{\tiny{[GM$\langle$P2]}}
                \texttt{VERSCHIEBE: S, (4,5)} \\
            }
        }
    }
     \\ \\

    \theutterance \stepcounter{utterance}  
    & & & \multicolumn{2}{p{0.3\linewidth}}{
        \cellcolor[rgb]{0.9,0.9,0.9}{
            \makecell[{{p{\linewidth}}}]{
                \texttt{\tiny{[GM$|$GM]}}
                \texttt{VERSCHIEBE: S, (4,5)} \\
            }
        }
    }
    & & \\ \\

    \theutterance \stepcounter{utterance}  
    & & & \multicolumn{2}{p{0.3\linewidth}}{
        \cellcolor[rgb]{0.9,0.9,0.9}{
            \makecell[{{p{\linewidth}}}]{
                \texttt{\tiny{[GM$|$GM]}}
                \texttt{'S' erfolgreich nach (4,5) verschoben. Dein Raster sieht jetzt so aus:} \\
\texttt{```} \\
\texttt{ 1234567} \\
\texttt{╔═╤═╤═╤═╗} \\
\texttt{║P│O│T│◌║ 1} \\
\texttt{╟─┴─┼─┤◌║ 2} \\
\texttt{║◌M◌│◌└─╢ 3} \\
\texttt{╟───┘◌◌◌║ 4} \\
\texttt{║◌◌◌S◌┌─╢ 5} \\
\texttt{║◌D◌◌A│◌║ 6} \\
\texttt{║◌◌◌◌◌│◌║ 7} \\
\texttt{╚═════╧═╝} \\
\texttt{```} \\
            }
        }
    }
    & & \\ \\

    \theutterance \stepcounter{utterance}  
    & & \multicolumn{4}{p{0.6\linewidth}}{
        \cellcolor[rgb]{0.9,0.9,0.9}{
            \makecell[{{p{\linewidth}}}]{
                \texttt{\tiny{[P1$\langle$GM]}}
                \texttt{Deine Nachricht wurde an deinen Mitspieler übermittelt.} \\
\\ 
\\ 
\texttt{Ihr seid gerade bei Runde 16 von maximal 28.} \\
\texttt{Ihr habt gemeinsam 6 von maximal 16 Strafpunkten angesammelt.} \\
\texttt{Dein Mitspieler hat ein Objekt auf seinem Raster bewegt.} \\
\\ 
\\ 
\texttt{Wie lautet dein nächster Befehl?} \\
            }
        }
    }
    & \\ \\

    \theutterance \stepcounter{utterance}  
    & \multicolumn{4}{p{0.6\linewidth}}{
        \cellcolor[rgb]{0.9,0.9,0.9}{
            \makecell[{{p{\linewidth}}}]{
                \texttt{\tiny{[P1$\rangle$GM]}}
                \texttt{VERSCHIEBE: T, (5,1)} \\
            }
        }
    }
    & & \\ \\

    \theutterance \stepcounter{utterance}  
    & & & \multicolumn{2}{p{0.3\linewidth}}{
        \cellcolor[rgb]{0.9,0.9,0.9}{
            \makecell[{{p{\linewidth}}}]{
                \texttt{\tiny{[GM$|$GM]}}
                \texttt{VERSCHIEBE: T, (5,1)} \\
            }
        }
    }
    & & \\ \\

    \theutterance \stepcounter{utterance}  
    & & & \multicolumn{2}{p{0.3\linewidth}}{
        \cellcolor[rgb]{0.9,0.9,0.9}{
            \makecell[{{p{\linewidth}}}]{
                \texttt{\tiny{[GM$|$GM]}}
                \texttt{'T' erfolgreich nach (5,1) verschoben. Dein Raster sieht jetzt so aus:} \\
\texttt{```} \\
\texttt{ 1234567} \\
\texttt{╔═╤═╤═╤═╗} \\
\texttt{║P│O│T│◌║ 1} \\
\texttt{╟─┴─┼─┤◌║ 2} \\
\texttt{║D◌M│A└─╢ 3} \\
\texttt{╟───┘◌S◌║ 4} \\
\texttt{║◌◌◌◌◌┌─╢ 5} \\
\texttt{║◌◌◌◌◌│◌║ 6} \\
\texttt{║◌◌◌◌◌│◌║ 7} \\
\texttt{╚═════╧═╝} \\
\texttt{```} \\
            }
        }
    }
    & & \\ \\

    \theutterance \stepcounter{utterance}  
    & & \multicolumn{4}{p{0.6\linewidth}}{
        \cellcolor[rgb]{0.9,0.9,0.9}{
            \makecell[{{p{\linewidth}}}]{
                \texttt{\tiny{[GM$\rangle$P2]}}
                \texttt{'S' erfolgreich nach (4,5) verschoben. Dein Raster sieht jetzt so aus:} \\
\texttt{```} \\
\texttt{ 1234567} \\
\texttt{╔═╤═╤═╤═╗} \\
\texttt{║P│O│T│◌║ 1} \\
\texttt{╟─┴─┼─┤◌║ 2} \\
\texttt{║◌M◌│◌└─╢ 3} \\
\texttt{╟───┘◌◌◌║ 4} \\
\texttt{║◌◌◌S◌┌─╢ 5} \\
\texttt{║◌D◌◌A│◌║ 6} \\
\texttt{║◌◌◌◌◌│◌║ 7} \\
\texttt{╚═════╧═╝} \\
\texttt{```} \\
\texttt{Ihr seid gerade bei Runde 17 von maximal 28.} \\
\texttt{Ihr habt gemeinsam 6 von maximal 16 Strafpunkten angesammelt.} \\
\texttt{Dein Mitspieler hat ein Objekt auf seinem Raster bewegt.} \\
\\ 
\\ 
\texttt{Wie lautet dein nächster Befehl?} \\
            }
        }
    }
    & \\ \\

    \theutterance \stepcounter{utterance}  
    & & & \multicolumn{4}{p{0.6\linewidth}}{
        \cellcolor[rgb]{0.9,0.9,0.9}{
            \makecell[{{p{\linewidth}}}]{
                \texttt{\tiny{[GM$\langle$P2]}}
                \texttt{VERSCHIEBE: D, (1,3)} \\
            }
        }
    }
     \\ \\

    \theutterance \stepcounter{utterance}  
    & & & \multicolumn{2}{p{0.3\linewidth}}{
        \cellcolor[rgb]{0.9,0.9,0.9}{
            \makecell[{{p{\linewidth}}}]{
                \texttt{\tiny{[GM$|$GM]}}
                \texttt{VERSCHIEBE: D, (1,3)} \\
            }
        }
    }
    & & \\ \\

    \theutterance \stepcounter{utterance}  
    & & & \multicolumn{2}{p{0.3\linewidth}}{
        \cellcolor[rgb]{0.9,0.9,0.9}{
            \makecell[{{p{\linewidth}}}]{
                \texttt{\tiny{[GM$|$GM]}}
                \texttt{'D' erfolgreich nach (1,3) verschoben. Dein Raster sieht jetzt so aus:} \\
\texttt{```} \\
\texttt{ 1234567} \\
\texttt{╔═╤═╤═╤═╗} \\
\texttt{║P│O│T│◌║ 1} \\
\texttt{╟─┴─┼─┤◌║ 2} \\
\texttt{║DM◌│◌└─╢ 3} \\
\texttt{╟───┘◌◌◌║ 4} \\
\texttt{║◌◌◌S◌┌─╢ 5} \\
\texttt{║◌◌◌◌A│◌║ 6} \\
\texttt{║◌◌◌◌◌│◌║ 7} \\
\texttt{╚═════╧═╝} \\
\texttt{```} \\
            }
        }
    }
    & & \\ \\

    \theutterance \stepcounter{utterance}  
    & & \multicolumn{4}{p{0.6\linewidth}}{
        \cellcolor[rgb]{0.9,0.9,0.9}{
            \makecell[{{p{\linewidth}}}]{
                \texttt{\tiny{[P1$\langle$GM]}}
                \texttt{'T' erfolgreich nach (5,1) verschoben. Dein Raster sieht jetzt so aus:} \\
\texttt{```} \\
\texttt{ 1234567} \\
\texttt{╔═╤═╤═╤═╗} \\
\texttt{║P│O│T│◌║ 1} \\
\texttt{╟─┴─┼─┤◌║ 2} \\
\texttt{║D◌M│A└─╢ 3} \\
\texttt{╟───┘◌S◌║ 4} \\
\texttt{║◌◌◌◌◌┌─╢ 5} \\
\texttt{║◌◌◌◌◌│◌║ 6} \\
\texttt{║◌◌◌◌◌│◌║ 7} \\
\texttt{╚═════╧═╝} \\
\texttt{```} \\
\texttt{Ihr seid gerade bei Runde 17 von maximal 28.} \\
\texttt{Ihr habt gemeinsam 6 von maximal 16 Strafpunkten angesammelt.} \\
\texttt{Dein Mitspieler hat ein Objekt auf seinem Raster bewegt.} \\
\\ 
\\ 
\texttt{Wie lautet dein nächster Befehl?} \\
            }
        }
    }
    & \\ \\

    \theutterance \stepcounter{utterance}  
    & \multicolumn{4}{p{0.6\linewidth}}{
        \cellcolor[rgb]{0.9,0.9,0.9}{
            \makecell[{{p{\linewidth}}}]{
                \texttt{\tiny{[P1$\rangle$GM]}}
                \texttt{VERSCHIEBE: S, (4,5)} \\
            }
        }
    }
    & & \\ \\

    \theutterance \stepcounter{utterance}  
    & & & \multicolumn{2}{p{0.3\linewidth}}{
        \cellcolor[rgb]{0.9,0.9,0.9}{
            \makecell[{{p{\linewidth}}}]{
                \texttt{\tiny{[GM$|$GM]}}
                \texttt{VERSCHIEBE: S, (4,5)} \\
            }
        }
    }
    & & \\ \\

    \theutterance \stepcounter{utterance}  
    & & & \multicolumn{2}{p{0.3\linewidth}}{
        \cellcolor[rgb]{0.9,0.9,0.9}{
            \makecell[{{p{\linewidth}}}]{
                \texttt{\tiny{[GM$|$GM]}}
                \texttt{'S' erfolgreich nach (4,5) verschoben. Dein Raster sieht jetzt so aus:} \\
\texttt{```} \\
\texttt{ 1234567} \\
\texttt{╔═╤═╤═╤═╗} \\
\texttt{║P│O│T│◌║ 1} \\
\texttt{╟─┴─┼─┤◌║ 2} \\
\texttt{║D◌M│A└─╢ 3} \\
\texttt{╟───┘◌◌◌║ 4} \\
\texttt{║◌◌◌S◌┌─╢ 5} \\
\texttt{║◌◌◌◌◌│◌║ 6} \\
\texttt{║◌◌◌◌◌│◌║ 7} \\
\texttt{╚═════╧═╝} \\
\texttt{```} \\
            }
        }
    }
    & & \\ \\

    \theutterance \stepcounter{utterance}  
    & & \multicolumn{4}{p{0.6\linewidth}}{
        \cellcolor[rgb]{0.9,0.9,0.9}{
            \makecell[{{p{\linewidth}}}]{
                \texttt{\tiny{[GM$\rangle$P2]}}
                \texttt{'D' erfolgreich nach (1,3) verschoben. Dein Raster sieht jetzt so aus:} \\
\texttt{```} \\
\texttt{ 1234567} \\
\texttt{╔═╤═╤═╤═╗} \\
\texttt{║P│O│T│◌║ 1} \\
\texttt{╟─┴─┼─┤◌║ 2} \\
\texttt{║DM◌│◌└─╢ 3} \\
\texttt{╟───┘◌◌◌║ 4} \\
\texttt{║◌◌◌S◌┌─╢ 5} \\
\texttt{║◌◌◌◌A│◌║ 6} \\
\texttt{║◌◌◌◌◌│◌║ 7} \\
\texttt{╚═════╧═╝} \\
\texttt{```} \\
\texttt{Ihr seid gerade bei Runde 18 von maximal 28.} \\
\texttt{Ihr habt gemeinsam 6 von maximal 16 Strafpunkten angesammelt.} \\
\texttt{Dein Mitspieler hat ein Objekt auf seinem Raster bewegt.} \\
\\ 
\\ 
\texttt{Wie lautet dein nächster Befehl?} \\
            }
        }
    }
    & \\ \\

    \theutterance \stepcounter{utterance}  
    & & & \multicolumn{4}{p{0.6\linewidth}}{
        \cellcolor[rgb]{0.9,0.9,0.9}{
            \makecell[{{p{\linewidth}}}]{
                \texttt{\tiny{[GM$\langle$P2]}}
                \texttt{VERSCHIEBE: A, (5,3)} \\
            }
        }
    }
     \\ \\

    \theutterance \stepcounter{utterance}  
    & & & \multicolumn{2}{p{0.3\linewidth}}{
        \cellcolor[rgb]{0.9,0.9,0.9}{
            \makecell[{{p{\linewidth}}}]{
                \texttt{\tiny{[GM$|$GM]}}
                \texttt{VERSCHIEBE: A, (5,3)} \\
            }
        }
    }
    & & \\ \\

    \theutterance \stepcounter{utterance}  
    & & & \multicolumn{2}{p{0.3\linewidth}}{
        \cellcolor[rgb]{0.9,0.9,0.9}{
            \makecell[{{p{\linewidth}}}]{
                \texttt{\tiny{[GM$|$GM]}}
                \texttt{'A' erfolgreich nach (5,3) verschoben. Dein Raster sieht jetzt so aus:} \\
\texttt{```} \\
\texttt{ 1234567} \\
\texttt{╔═╤═╤═╤═╗} \\
\texttt{║P│O│T│◌║ 1} \\
\texttt{╟─┴─┼─┤◌║ 2} \\
\texttt{║DM◌│A└─╢ 3} \\
\texttt{╟───┘◌◌◌║ 4} \\
\texttt{║◌◌◌S◌┌─╢ 5} \\
\texttt{║◌◌◌◌◌│◌║ 6} \\
\texttt{║◌◌◌◌◌│◌║ 7} \\
\texttt{╚═════╧═╝} \\
\texttt{```} \\
            }
        }
    }
    & & \\ \\

    \theutterance \stepcounter{utterance}  
    & & \multicolumn{4}{p{0.6\linewidth}}{
        \cellcolor[rgb]{0.9,0.9,0.9}{
            \makecell[{{p{\linewidth}}}]{
                \texttt{\tiny{[P1$\langle$GM]}}
                \texttt{'S' erfolgreich nach (4,5) verschoben. Dein Raster sieht jetzt so aus:} \\
\texttt{```} \\
\texttt{ 1234567} \\
\texttt{╔═╤═╤═╤═╗} \\
\texttt{║P│O│T│◌║ 1} \\
\texttt{╟─┴─┼─┤◌║ 2} \\
\texttt{║D◌M│A└─╢ 3} \\
\texttt{╟───┘◌◌◌║ 4} \\
\texttt{║◌◌◌S◌┌─╢ 5} \\
\texttt{║◌◌◌◌◌│◌║ 6} \\
\texttt{║◌◌◌◌◌│◌║ 7} \\
\texttt{╚═════╧═╝} \\
\texttt{```} \\
\texttt{Ihr seid gerade bei Runde 18 von maximal 28.} \\
\texttt{Ihr habt gemeinsam 6 von maximal 16 Strafpunkten angesammelt.} \\
\texttt{Dein Mitspieler hat ein Objekt auf seinem Raster bewegt.} \\
\\ 
\\ 
\texttt{Wie lautet dein nächster Befehl?} \\
            }
        }
    }
    & \\ \\

    \theutterance \stepcounter{utterance}  
    & \multicolumn{4}{p{0.6\linewidth}}{
        \cellcolor[rgb]{0.9,0.9,0.9}{
            \makecell[{{p{\linewidth}}}]{
                \texttt{\tiny{[P1$\rangle$GM]}}
                \texttt{SPRICH: Ich habe alle Objekte verschoben: P(1,1), O(3,1), T(5,1), S(4,5), D(1,3), A(5,3), M(3,3). Bist du auch fertig? Wenn ja, antworte mit SPRICH: fertig!} \\
            }
        }
    }
    & & \\ \\

    \theutterance \stepcounter{utterance}  
    & & & \multicolumn{2}{p{0.3\linewidth}}{
        \cellcolor[rgb]{0.9,0.9,0.9}{
            \makecell[{{p{\linewidth}}}]{
                \texttt{\tiny{[GM$|$GM]}}
                \texttt{SPRICH: Ich habe alle Objekte verschoben: P(1,1), O(3,1), T(5,1), S(4,5), D(1,3), A(5,3), M(3,3). Bist du auch fertig? Wenn ja, antworte mit SPRICH: fertig!} \\
            }
        }
    }
    & & \\ \\

    \theutterance \stepcounter{utterance}  
    & & & \multicolumn{2}{p{0.3\linewidth}}{
        \cellcolor[rgb]{0.9,0.9,0.9}{
            \makecell[{{p{\linewidth}}}]{
                \texttt{\tiny{[GM$|$GM]}}
                \texttt{Invalid format: head is not empty: 'SPRICH: Ich habe alle Objekte verschoben: P(1,1), O(3,1), T(5,1), S(4,5), D(1,3), A(5,3), M(3,3). Bist du auch fertig? Wenn ja, antworte mit '} \\
            }
        }
    }
    & & \\ \\

    \theutterance \stepcounter{utterance}  
    & & \multicolumn{4}{p{0.6\linewidth}}{
        \cellcolor[rgb]{0.9,0.9,0.9}{
            \makecell[{{p{\linewidth}}}]{
                \texttt{\tiny{[P1$\langle$GM]}}
                \texttt{Strafpunkt: Deine Nachricht darf nichts vor dem Befehl enthalten! Stelle sicher, dass deine Antwort ausschließlich einen der beiden Befehle `SPRICH: <NACHRICHT>` und `VERSCHIEBE: <OBJECT>, (<X>, <Y>)` und sonst nichts enthält!} \\
\\ 
\texttt{Ihr habt gemeinsam 7 von maximal 16 Strafpunkten angesammelt.} \\
\texttt{Bitte versuche es noch mal!} \\
            }
        }
    }
    & \\ \\

    \theutterance \stepcounter{utterance}  
    & \multicolumn{4}{p{0.6\linewidth}}{
        \cellcolor[rgb]{0.9,0.9,0.9}{
            \makecell[{{p{\linewidth}}}]{
                \texttt{\tiny{[P1$\rangle$GM]}}
                \texttt{SPRICH: fertig?} \\
            }
        }
    }
    & & \\ \\

    \theutterance \stepcounter{utterance}  
    & & & \multicolumn{2}{p{0.3\linewidth}}{
        \cellcolor[rgb]{0.9,0.9,0.9}{
            \makecell[{{p{\linewidth}}}]{
                \texttt{\tiny{[GM$|$GM]}}
                \texttt{SPRICH: fertig?} \\
            }
        }
    }
    & & \\ \\

    \theutterance \stepcounter{utterance}  
    & & \multicolumn{4}{p{0.6\linewidth}}{
        \cellcolor[rgb]{0.9,0.9,0.9}{
            \makecell[{{p{\linewidth}}}]{
                \texttt{\tiny{[GM$\rangle$P2]}}
                \texttt{'A' erfolgreich nach (5,3) verschoben. Dein Raster sieht jetzt so aus:} \\
\texttt{```} \\
\texttt{ 1234567} \\
\texttt{╔═╤═╤═╤═╗} \\
\texttt{║P│O│T│◌║ 1} \\
\texttt{╟─┴─┼─┤◌║ 2} \\
\texttt{║DM◌│A└─╢ 3} \\
\texttt{╟───┘◌◌◌║ 4} \\
\texttt{║◌◌◌S◌┌─╢ 5} \\
\texttt{║◌◌◌◌◌│◌║ 6} \\
\texttt{║◌◌◌◌◌│◌║ 7} \\
\texttt{╚═════╧═╝} \\
\texttt{```} \\
\texttt{Ihr seid gerade bei Runde 19 von maximal 28.} \\
\texttt{Ihr habt gemeinsam 7 von maximal 16 Strafpunkten angesammelt.} \\
\texttt{Dein Mitspieler hat dir diese Nachricht geschickt:} \\
\\ 
\\ 
\texttt{"fertig?"} \\
\\ 
\\ 
\texttt{Wie lautet dein nächster Befehl?} \\
            }
        }
    }
    & \\ \\

    \theutterance \stepcounter{utterance}  
    & & & \multicolumn{4}{p{0.6\linewidth}}{
        \cellcolor[rgb]{0.9,0.9,0.9}{
            \makecell[{{p{\linewidth}}}]{
                \texttt{\tiny{[GM$\langle$P2]}}
                \texttt{SPRICH: fertig!} \\
            }
        }
    }
     \\ \\

    \theutterance \stepcounter{utterance}  
    & & & \multicolumn{2}{p{0.3\linewidth}}{
        \cellcolor[rgb]{0.9,0.9,0.9}{
            \makecell[{{p{\linewidth}}}]{
                \texttt{\tiny{[GM$|$GM]}}
                \texttt{SPRICH: fertig!} \\
            }
        }
    }
    & & \\ \\

    \theutterance \stepcounter{utterance}  
    & & & \multicolumn{2}{p{0.3\linewidth}}{
        \cellcolor[rgb]{0.9,0.9,0.9}{
            \makecell[{{p{\linewidth}}}]{
                \texttt{\tiny{[GM$|$GM]}}
                \texttt{true} \\
            }
        }
    }
    & & \\ \\

    \theutterance \stepcounter{utterance}  
    & & & \multicolumn{2}{p{0.3\linewidth}}{
        \cellcolor[rgb]{0.9,0.9,0.9}{
            \makecell[{{p{\linewidth}}}]{
                \texttt{\tiny{[GM$|$GM]}}
                \texttt{* success: True} \\
\texttt{* lose: False} \\
\texttt{* aborted: False} \\
\texttt{{-}{-}{-}{-}{-}{-}{-}} \\
\texttt{* Shifts: 11.00} \\
\texttt{* Max Shifts: 12.00} \\
\texttt{* Min Shifts: 6.00} \\
\texttt{* End Distance Sum: 1.00} \\
\texttt{* Init Distance Sum: 18.94} \\
\texttt{* Expected Distance Sum: 29.33} \\
\texttt{* Penalties: 7.00} \\
\texttt{* Max Penalties: 16.00} \\
\texttt{* Rounds: 19.00} \\
\texttt{* Max Rounds: 28.00} \\
\texttt{* Object Count: 7.00} \\
            }
        }
    }
    & & \\ \\

    \theutterance \stepcounter{utterance}  
    & & & \multicolumn{2}{p{0.3\linewidth}}{
        \cellcolor[rgb]{0.9,0.9,0.9}{
            \makecell[{{p{\linewidth}}}]{
                \texttt{\tiny{[GM$|$GM]}}
                \texttt{Initial Grids:} \\
\texttt{```} \\
\texttt{Player 1:     Player 2:} \\
\texttt{ 1234567       1234567} \\
\texttt{╔═╤═╤═╤═╗     ╔═╤═╤═╤═╗} \\
\texttt{║◌│◌│◌│A║ 1   ║◌│◌│◌│◌║ 1} \\
\texttt{╟─┴─┼─┤◌║ 2   ╟─┴─┼─┤◌║ 2} \\
\texttt{║◌◌M│◌└─╢ 3   ║◌M◌│◌└─╢ 3} \\
\texttt{╟───┘◌S◌║ 4   ╟───┘P◌◌║ 4} \\
\texttt{║◌◌◌◌◌┌─╢ 5   ║O◌◌◌◌┌─╢ 5} \\
\texttt{║◌◌◌PT│O║ 6   ║◌D◌◌A│S║ 6} \\
\texttt{║◌D◌◌◌│◌║ 7   ║◌◌◌◌T│◌║ 7} \\
\texttt{╚═════╧═╝     ╚═════╧═╝} \\
\texttt{```} \\
\\ 
\texttt{Final Grids:} \\
\texttt{```} \\
\texttt{Player 1:     Player 2:} \\
\texttt{ 1234567       1234567} \\
\texttt{╔═╤═╤═╤═╗     ╔═╤═╤═╤═╗} \\
\texttt{║P│O│T│◌║ 1   ║P│O│T│◌║ 1} \\
\texttt{╟─┴─┼─┤◌║ 2   ╟─┴─┼─┤◌║ 2} \\
\texttt{║D◌M│A└─╢ 3   ║DM◌│A└─╢ 3} \\
\texttt{╟───┘◌◌◌║ 4   ╟───┘◌◌◌║ 4} \\
\texttt{║◌◌◌S◌┌─╢ 5   ║◌◌◌S◌┌─╢ 5} \\
\texttt{║◌◌◌◌◌│◌║ 6   ║◌◌◌◌◌│◌║ 6} \\
\texttt{║◌◌◌◌◌│◌║ 7   ║◌◌◌◌◌│◌║ 7} \\
\texttt{╚═════╧═╝     ╚═════╧═╝} \\
\texttt{```} \\
            }
        }
    }
    & & \\ \\

\end{supertabular}
}

\end{document}
